\section{Classical inputs and sphere normalizations}

The framework above fixes the classical and synthetic categories, the
maps and detecting classes, and the realization functor. We now state
the classical inputs and the further comparison assumptions used with
those constructions. Statements quoted from \cite{LWX126} use its
$2$-completed, connective, finite-type hypotheses; different domains
are stated at the individual entries. Completeness enters the passage
to untruncated lifts, while finite-quotient formulas have their own
stated ranges.

Write $E_r^{s,t}(X)$ for the classical mod-two Adams page, with stem
$t-s$, and $\pi_{m,w}U$ for synthetic homotopy. The synthetic analogue
$\nu X$ and the Hopf map $\nu:S^3\to S^0$ are distinguished by their
arguments. The class called $\tau$ in \cite{BHS} is denoted $\lambda$
here. Paper numbering refers to \cite{LWX126}; see the
\href{paper/2412.10879.pdf}{original paper}.

The main sphere argument uses the Hopf and theta-five classes, the tmf
calculations, the two infinite-page window comparisons
(Propositions~\ref{sem:Literature:infinite-nu-window}
and~\ref{sem:Literature:infinite-finite-window}), the lifting theorem
\ref{sem:Literature:finite-bockstein}, and the survival criterion
\ref{sem:Literature:bx-survival}. The geometric bridge is stated before
the finite second-line inputs so that the exclusion of $h_7^2$ uses
results already introduced.

The Hopf classes and theta-five supply the homotopy classes used in the
obstruction. We first recall their classical detections and normalize
their synthetic lifts. The vanishing range and the one-line
differentials restrict the other sphere classes that can occur.

\begin{theorem}[Connectivity and mod-two homology]

\label{sem:Literature:connected-homology}
\level{coarse}
\uses{sem:AdamsSynthetic:categories}
The connectivity statement follows from \cite[Proposition II.5.22]{SchwedeSymmetricSpectra}; the cohomology comparison is the field case of \cite[Theorem IV.4.5]{EKMM}.\par

The sphere is $(-1)$-connected. If $X$ is $n$-connected, then
$H_i(X;\mathbb F_2)=0$ for $i\le n$. There is an additive
universal-coefficient equivalence
\[
 [\Sigma^nX,H\mathbb F_2]\cong
 \operatorname{Hom}(H_{-n}(X;\mathbb F_2),\mathbb F_2).
\]
We use conventional cohomological indexing
$H^q(X;\mathbb F_2)=[X,\Sigma^qH\mathbb F_2]$ throughout.
Thus the group on the left is $H^{-n}(X;\mathbb F_2)$, explaining
the homology index $-n$. We use the same indexing for the Steenrod
modules below.


\lean{KIP126.Literature.Standard.sphere_connected,KIP126.Literature.Standard.connected_homology_vanishing,KIP126.Literature.Standard.universal_coefficient}
\end{theorem}
\begin{proof}
\uses{sem:AdamsSynthetic:paper-spectrum-instances}
The sphere has no negative homotopy groups. The connectivity statement
for homology follows from the connective coefficient spectrum
$H\mathbb F_2$. Over the field $\mathbb F_2$, the universal-coefficient
comparison identifies cohomology with the dual of homology in the
displayed degree.
\end{proof}

\begin{proposition}[Low-weight sphere comparison isomorphisms]
\label{sem:Literature:low-sphere-comparisons}
\level{coarse}
\uses{sem:AdamsSynthetic:nu,paper:3.2}
These are the unit and suspension comparisons of $\nu$; see
\cite[Definition 3.2]{LWX126}.\par

For $k=1,2,3$, the unit and suspension comparisons give isomorphisms
\[\nu(S^k)\cong S^{k,k}.\]
We use these three cases, with the sphere notation of
Definition~\ref{paper:3.2}.

\lean{KIP126.Literature.nuSphereOne,KIP126.Literature.nuSphereTwo,KIP126.Literature.nuSphereThree}
\end{proposition}

\begin{definition}[Realization in sphere coordinates]
\label{sem:Literature:realization-coordinates}
\level{coarse}
\uses{sem:AdamsSynthetic:recovery,paper:3.2,paper:3.3}
See \cite[Proposition 4.40]{Pst}.\par

Let $\operatorname{Re}$ be the additive realization functor. Assume sphere
isomorphisms
$\rho_{m,w}:S^m\xrightarrow{\sim}\operatorname{Re}(S^{m,w})$
under which every $\lambda^k$ realizes to the identity. At the unit,
require agreement with the zero-shift, inverse
$\nu$--realization and synthetic-unit comparisons. Write
$\operatorname{Re}(a)$ for the resulting realization of a sphere class
or a class on $\nu X$. Put $S^m_{\mathrm{syn}}=\Sigma^m\mathbf1$,
the ordinary suspension by $m$ of the synthetic tensor unit. If
$c_m:S^{m,0}\xrightarrow{\sim}S^m_{\mathrm{syn}}$
is the comparison with ordinary synthetic suspension, define
the ordinary sphere coordinate by
\[
r_m=\operatorname{Re}(c_m)\circ\rho_{m,0}:
S^m\xrightarrow{\sim}\operatorname{Re}(S^m_{\mathrm{syn}}).
\]

\lean{KIP126.Literature.Standard.realization_additive,KIP126.Literature.Standard.realization_sphere,KIP126.Literature.Standard.realization_unit,KIP126.Literature.Standard.realization_lambda,KIP126.Literature.Standard.realizationData,KIP126.Literature.Standard.realizeSphere,KIP126.Literature.Standard.realizeNu,KIP126.Literature.realizationSphereCoordinate}
\end{definition}


\begin{proposition}[The Hopf classes and their synthetic normalization]
\label{sem:Literature:hopf-classes}
\level{coarse}
\uses{sem:AdamsSynthetic:standard-classes,sem:AdamsSynthetic:normalized-map,sem:Literature:realization-coordinates,sem:AdamsSynthetic:e2-comparison}
For the classical Hopf classes and their synthetic lifts, see
\cite[Lemma 2.3]{Adams58} and \cite[Proposition 3.16]{LWX126}.\par

The classical maps $\eta$ and $\nu$ are detected by the nonzero permanent cycles $h_1\in E_2^{1,2}$ and $h_2\in E_2^{1,4}$, respectively. The normalized synthetic lift $\eta\in\pi_{1,2}S^{0,0}$ has filtration exponent one, first-quotient label $h_1$, and realization equal to the classical $\eta$.

The first-quotient and realization equalities are part of the assumed
normalization of this synthetic lift.

\lean{KIP126.Literature.Standard.eta,KIP126.Literature.Standard.eta_normalized,KIP126.Literature.Standard.eta_detection,KIP126.Literature.Standard.eta_exponent,KIP126.Literature.Standard.eta_permanent,KIP126.Literature.Standard.eta_label,KIP126.Literature.Standard.eta_realization,KIP126.Literature.Standard.nu_detection,KIP126.Literature.Standard.nu_permanent}
\end{proposition}

\begin{proposition}[The class theta-five and the filtration of its choices]
\label{sem:Literature:theta-five}
\level{coarse}
\uses{sem:AdamsSynthetic:standard-classes,sem:AdamsSynthetic:selected-convergence,sem:Literature:realization-coordinates,sem:AdamsSynthetic:e2-comparison}
The classical existence theorem is \cite[Theorem 2.1]{BJMtheta5}, with
the detection in \cite[Corollary 1.3]{Xu}. We use the synthetic lift
under the comparisons above.\par

There is a classical class $\theta_5\in\pi_{62}S^0$ of order two detected by the nonzero permanent cycle $h_5^2\in E_2^{2,64}$. Fix such a class. There is a synthetic lift $a\in\pi_{62,64}S^{0,0}$ with first-quotient label $h_5^2$ and $\operatorname{Re}(a)=\theta_5$.

For the two-primary stem calculation, see \cite[Corollary 1.3 and Table 1]{IWX} and \cite[Lemma 7.10, proof]{LWX126}.\par

Every class $a\in\pi_{62}S^0$ annihilated by a power of two satisfies $2a=0$. If $a,b\in\pi_{62}S^0$ are detected by $h_5^2$, then $a-b$ has Adams filtration at least six.


\lean{KIP126.Literature.Standard.theta5,KIP126.Literature.Standard.h5Square_permanent,KIP126.Literature.Standard.theta5_detection,KIP126.Literature.Standard.theta5_order_two,KIP126.Literature.Standard.theta5_lift,KIP126.Literature.Standard.theta5_detected_exists,KIP126.Literature.Standard.stem62_exponent_two,KIP126.Literature.Standard.theta5_filtration_gap,KIP126.Literature.Standard.theta5Construction}
\end{proposition}
\begin{proof}
\uses{sem:AdamsSynthetic:e2-comparison,sem:Literature:realization-coordinates}
The existence, detection and order of $\theta_5$ are the BJM--Xu result.
The stem calculation gives the two-primary exponent and the filtration
gap between two classes with detector $h_5^2$.
The synthetic lift uses the first-quotient and realization comparisons.
The exponent assertion applies to the two-primary subgroup of the
integral stem.
\end{proof}

\begin{theorem}[The positive-stem Adams vanishing range]
\label{sem:Literature:sphere-vanishing}
\level{coarse}
\uses{sem:AdamsSynthetic:families}
Adams' vanishing line is recorded in \cite[Proposition 3.3(a)]{MathewTorsion}.\par

In the classical sphere Adams spectral sequence,
\[0<t-s<2s-3\quad\Longrightarrow\quad E_2^{s,t}(S^0)=0.\]


\lean{KIP126.Literature.Standard.sphereVanishing}
\end{theorem}

\begin{theorem}[The Hopf-invariant-one differentials]
\label{sem:Literature:one-line-differentials}
\level{coarse}
\uses{sem:AdamsSynthetic:standard-classes,sem:AdamsSynthetic:selected-convergence}
The degree-two detection is \cite[Lemma 2.3]{Adams58}. The differential family is due to \cite{Adams60}; see \cite[Theorem 1.2]{BurklundXu}.\par

The class $h_0$ detects the classical degree-two map. For each integer $4\le j\le7$, there is a nonzero differential
\[d_2(h_j)=h_0h_{j-1}^2\in E_2^{3,1+2^j}(S^0).\]


\lean{KIP126.Literature.Standard.two_detection,KIP126.Literature.Standard.adamsOneLine_d2}
\end{theorem}

\section{The tmf detector}
The tmf Hurewicz map tests two possible sphere phenomena. Vanishing in
stem $62$ controls the image of theta-five, while a nonzero image in
stem $125$ detects a class that might otherwise be hidden in high Adams
filtration. We first identify the detector with tmf, including its unit
map.

\begin{definition}[The tmf detector and its Adams tower]
\label{sem:Literature:tmf-unit}
\label{sem:Literature:tmf-convergence}
\level{coarse}
\uses{sem:AdamsSynthetic:categories,sem:AdamsSynthetic:families,sem:AdamsSynthetic:selected-convergence,sem:AdamsSynthetic:paper-spectrum-instances}
See \cite[Sections 0.1 and 1.2]{BR21}. Let $D$ denote the detector
spectrum fixed above. Assume a unital associative algebra spectrum
$\mathit{tmf}$ identified, together with its unit
$u:S^0\to\mathit{tmf}$, with topological modular forms. Assume also
an isomorphism $e:\mathit{tmf}\xrightarrow{\sim}D$ that sends $u$ to
the detector unit.

This spectrum is connective, has finite mod-two homology in each
degree, and its classical Adams tower converges to its homotopy groups
with the Adams filtration. Transport along $e$ preserves tower lifts
and associated-graded representatives.

All homotopy and Adams comparisons below use these identifications.

\lean{KIP126.Literature.Standard.tmf,KIP126.Literature.Standard.tmfIso,KIP126.Literature.Standard.tmf_algebra,KIP126.Literature.Standard.tmfUnit,KIP126.Literature.Standard.tmfUnit_agrees,KIP126.Literature.Standard.tmf_algebra_unit,KIP126.Literature.Standard.tmfObjectConstruction,KIP126.Literature.Standard.tmf_connective,KIP126.Literature.Standard.tmf_finiteModTwoType,KIP126.Literature.Standard.tmf_convergence}
\end{definition}
\begin{proof}
Connectivity, finite type and convergence transfer from $D$ along $e$.
The isomorphism $e$, and its tensor with $H\mathbb F_2$, identify the
homotopy and coefficient homology groups. The paper-range properties of
$D$ therefore give connectivity and finite type. Transport its tower
convergence along the same isomorphism; the induced tower maps preserve
the lifts and cofiber representatives of associated-graded classes.
\end{proof}

\begin{proposition}[The classical classes detected by g and Delta h-one g]
\label{sem:Literature:kappa-detection}
\level{coarse}
\uses{sem:AdamsSynthetic:families,sem:AdamsSynthetic:selected-convergence}
The detection of $\bar\kappa$ is given in \cite[Theorem 11.61(20)]{BR21}.\par

The group $E_2^{4,24}(S^0)$ has a unique nonzero element $g$. It detects a class $\bar\kappa\in\pi_{20}S^0$, which we fix.

For the stem-$45$ class and its detection, see \cite[Theorem 11.59 and Tables 1.1 and 11.3]{BR21}.\par

The group $E_2^{9,54}(S^0)$ has a unique nonzero element $\Delta h_1g$. It detects a class $w\in\pi_{45}S^0$, which we fix.


\lean{KIP126.Literature.Standard.g,KIP126.Literature.Standard.g_nonzero,KIP126.Literature.Standard.g_unique,KIP126.Literature.Standard.kappaBar,KIP126.Literature.Standard.kappaBar_detection,KIP126.Literature.Standard.deltaH1g,KIP126.Literature.Standard.deltaH1g_nonzero,KIP126.Literature.Standard.deltaH1g_unique,KIP126.Literature.Standard.wClass,KIP126.Literature.Standard.w_detection,KIP126.Literature.Standard.gConstruction,KIP126.Literature.Standard.deltaH1gConstruction,KIP126.Literature.Standard.kappaBarConstruction,KIP126.Literature.Standard.wConstruction}
\end{proposition}
\begin{proof}
\uses{sem:AdamsSynthetic:families,sem:AdamsSynthetic:selected-convergence}
The cited calculations identify the unique nonzero elements in the two
Adams groups and give their homotopy detections. We use the resulting
representatives $\bar\kappa$ and $w$ in these bidegrees.
\end{proof}

\begin{proposition}[The tmf Hurewicz input in stems 62 and 125]
\label{sem:Literature:tmf-hurewicz}
\level{coarse}
\uses{sem:Literature:tmf-unit,sem:Literature:tmf-convergence,sem:Literature:kappa-detection}
The homotopy calculation is \cite[Figure 1.1]{tmf}; the stem-$125$ product is identified in \cite[Theorem 11.80 and Proposition 11.83]{BR21}.\par

For tmf with the identification of
Definition~\ref{sem:Literature:tmf-unit}, assume the following
homotopy and Adams-page assertions.
One has $\pi_{62}\mathit{tmf}=0$ and $E_2^{s,63+s}(\mathit{tmf})=0$ for $s\le0$. Let $v=\bar\kappa^4w\in\pi_{125}S^0$. For every positive odd integer $n$,
\[n\,u_*(v)\ne0\quad\text{in }\pi_{125}\mathit{tmf}.\]

The stem-$62$ vanishing is read from \cite[Figure 1.1]{tmf}.
For the Adams assertion, \cite[Section 0.5]{BR21} identifies the
second page with $\operatorname{Ext}_{A(2)}(\mathbb F_2,\mathbb F_2)$;
negative Ext degrees vanish and $\operatorname{Ext}_{A(2)}^{0,63}=0$.
The nonzero stem-$125$ product is the class in
\cite[Proposition 11.83]{BR21}. Odd multiplication is an automorphism
of its finite two-primary homotopy group.

\lean{KIP126.Literature.Standard.tmf_vanishing62,KIP126.Literature.Standard.tmf_low_filtration63,KIP126.Literature.Standard.high125,KIP126.Literature.Standard.tmf_high125_nonzero}
\end{proposition}

\begin{theorem}[A third-page differential and its detector image]
\label{sem:Literature:tmf-d-three}
\level{coarse}
\uses{sem:Literature:tmf-unit,sem:AdamsSynthetic:families}
The calculation is \cite[Theorem 5.18 and Table 5.4]{BR21}.\par

Assume that there are classes
\[
 a\in E_2^{16,112}(\mathit{tmf}),\qquad
 b\in E_2^{19,114}(\mathit{tmf}),
\]
with third-page representatives $a_3,b_3$ satisfying $d_3(a_3)=b_3$.
Thus, in second-page label notation, $d_3(a)=b$. Under the
isomorphism $e:\mathit{tmf}\xrightarrow{\sim}D$, the labels $e_2(a)$
and $e_2(b)$ satisfy the same third-page differential relation.
Assume that the Adams comparisons identify $a,b$ with the classes
$w_2^2,\beta g^4$ in the cited calculation
$d_3(w_2^2)=\beta g^4$.


\lean{KIP126.Literature.Standard.tmf_v2_sixteen_d3,KIP126.Literature.Standard.detector_d3,KIP126.Literature.Standard.tmfV2Sixteen,KIP126.Literature.Standard.tmfBetaGFour,KIP126.Literature.Standard.tmfDifferentialConstruction}
\end{theorem}
\begin{proof}
\uses{sem:Literature:tmf-unit,sem:AdamsSynthetic:families}
Apply the map $e_3$ to the assumed equation:
$d_3(e_3a_3)=e_3(d_3a_3)=e_3b_3$. Naturality of the Adams spectral
sequence identifies their second-page labels with $e_2(a)$ and $e_2(b)$.
This is a relation with the specified source and target representatives;
it also allows zero representatives.
\end{proof}

\section{Synthetic Adams pages and finite lambda quotients}
The page comparisons of \cite{BHS}, Theorem A.8 and Corollaries A.9--A.11,
translate classical labels and differentials into synthetic weights.
The comparisons keep track of the classical detecting classes as well
as the page groups. We first describe finite pages, then the infinite-page
windows and their maps, and finally the finite lifts detected by a
Bockstein boundary. Completeness and the passage to untruncated homotopy
are treated in the next section.

\subsection{Finite pages and their comparison ranges}

\begin{definition}[Steenrod modules and Ext]
\label{sem:Literature:steenrod-ext}
\level{coarse}
\uses{sem:AdamsSynthetic:paper-spectrum-scope}
Let $H$ be the mod-two Eilenberg--Mac Lane spectrum in the completed
stable category. The category of stable mod-two operations has objects
$\Sigma^nH$, for $n\in\mathbb Z$, and stable maps as morphisms.
A Steenrod module is an additive functor from this operation category to
abelian groups. The cohomology module of $X$ takes the value
$[X,\Sigma^nH]$ in degree $n$, with operations acting by postcomposition.
These additive modules are closed under pointwise zero objects, kernels,
cokernels and finite products, giving the abelian category in which we
define
\[
\operatorname{Ext}_A^{s,t}(X)
 =\operatorname{Ext}^s\bigl(H^*(X;\mathbb F_2),
     H^*(\Sigma^tS^0;\mathbb F_2)\bigr)
\]
in this additive module category, using derived Ext also for integer $s$.
Thus the internal degree is determined by the suspension of the sphere.


\lean{KIP126.Literature.SteenrodExt.additiveModulesClosed,KIP126.Literature.SteenrodExt.modulesAbelian,KIP126.Literature.SteenrodExt.ext}
\end{definition}

\begin{proposition}[Synthetic Adams coordinates and their weight shifts]
\label{sem:Literature:nu-page-coordinates}
\level{coarse}
\uses{sem:AdamsSynthetic:families,sem:AdamsSynthetic:e2-comparison,sem:AdamsSynthetic:shift-coherence,sem:AdamsSynthetic:ess-shifts}
The synthetic coordinates are those of \cite[Theorem A.8]{BHS}.\par

For $X$ in the bounded-below finite-type spectrum range, use the
comparisons $K_{X,a,k}$ of
Definition~\ref{sem:AdamsSynthetic:e2-comparison}. For
$x\in E_2^{s,t}(X)$ and $k\ge0$, write $\lambda^kx$ for its label
$K_{X,0,k}(x)\in E_2^{s,t,t-k}(\nu X)$.
Weight shift induces isomorphisms
\[E_r^{s,t,w}(\nu X)\cong E_r^{s,t,w-1}(\Sigma^{0,-1}\nu X).\]
In these coordinates, the map induced by
$\lambda:\Sigma^{0,-1}\nu X\to\nu X$ lowers weight by one.

For $s,t\in\mathbb Z$ and $k\in\mathbb N$, this shift carries
the label of $x$ in degree $(s,t,t-k)$ to its label on
$\Sigma^{0,-1}\nu X$ in degree $(s,t,t-k-1)$, as in
\cite[Theorem A.8(2)]{BHS}.


\lean{KIP126.Literature.Standard.nuPage_weightShift,KIP126.Literature.Standard.nuPageData,KIP126.Literature.Standard.pageLambda,KIP126.Literature.Standard.classicalDifferential,KIP126.Literature.Standard.syntheticDifferential,KIP126.Literature.Standard.nuPage_e2}
\end{proposition}

% AimPaper: 3.6 theorem 799 807
\begin{theorem}[Synthetic Adams pages and differentials]
\label{paper:3.6}
\source{AimPaper Theorem 3.6}
\level{coarse}
For the page formulas, see \cite[Theorem A.8]{BHS} and
\cite[Theorem 3.6]{LWX126}
(\href{paper/main.tex.html#L799}{original statement}).\par
\label{thm:rigid}
    For every $2$-completed connective spectrum $X$ of degreewise finite
    mod-two type, the synthetic Adams spectral sequence for $\nu X$
     $$\SynSS_2^{s,t,w}(\nu X)\Longrightarrow \pi_{t-s,w}(\nu X)$$
     has an $E_2$-page
    $$\SynSS_2^{*,*,*}(\nu X) \cong E_2^{*,*}(X)\otimes \mathbb{F}_2[\lambda]$$
    with differentials
    $$d_r: \SynSS_r^{s,t,w}\to \SynSS_r^{s+r,t+r-1,w}.$$
    Here $E_2^{*,*}(X) \cong \Ext_A^{*,*}(X)$ is the classical Adams $E_2$-page of $X$. An element in $E_2^{s,t}(X)$ is viewed as having tridegree $(s, t, t)$ in $\SynSS_2^{s,t,t}(\nu X)$, and $\lambda$ is in tridegree $(0, 0, -1)$. Given a classical Adams differential $d_r^{\mathrm{cl}}(x)=y$, the corresponding synthetic differential is $d_r(x)=\lambda^{r-1}y$, which is $\lambda$-linear, and all synthetic Adams differentials arise in this way.

\uses{paper:3.3,sem:AdamsSynthetic:paper-spectrum-scope,sem:AdamsSynthetic:families,sem:AdamsSynthetic:e2-comparison,sem:Literature:nu-page-coordinates,sem:Literature:steenrod-ext,sem:AdamsSynthetic:selected-convergence}

The abutment uses the synthetic Adams filtration of
$\pi_{t-s,w}(\nu X)$: this filtration is exhaustive, separated and
complete, and its associated graded is canonically the $E_\infty$-page.

The $\lambda$ comparisons preserve representatives on $E_2$ and on
every later finite page, and the differential description applies to
every page element. The page and differential formulas are those of
\cite{BHS}, Theorem A.8; the strong-convergence comparison is
Proposition A.16, under the completion and convergence hypotheses
stated here.

\lean{KIP126.Literature.Standard.synthetic_adams_theorem}
\end{theorem}
\begin{proof}
Apply \cite{BHS}, Theorem A.8, to the $H\mathbb F_2$-synthetic analogue of $X$. Translate the BHS coordinates $(s,k,\omega)$ to the present coordinates by $k=t-s$ and $\omega=w-(t-s)$, and rename $\tau$ as $\lambda$. The free polynomial description then places a classical class in tridegree $(s,t,t)$ and $\lambda$ in tridegree $(0,0,-1)$.

The BHS differential degree becomes $(r,r-1,0)$. Consequently a classical differential $d_r(x)=y$ becomes $d_r(x)=\lambda^{r-1}y$: the indicated power is exactly what preserves weight. The theorem also gives $\lambda$-linearity and exhausts the synthetic differentials. For the abutment, apply the strong-convergence comparison with the
completion and convergence hypotheses in the statement.
\end{proof}



\begin{proposition}[Lambda inversion and the free comparison range]
\label{sem:Literature:lambda-inversion}
\label{sem:Literature:free-lambda-range}
\level{coarse}
\uses{sem:AdamsSynthetic:recovery,sem:Literature:nu-page-coordinates}
The localization comparison is \cite[Proposition 4.40]{Pst}; its page formulas follow from \cite[Theorem A.8 and Corollary A.9]{BHS}.
Inverting $\lambda$ compares the synthetic and classical spectral
sequences through maps of their cycle and boundary groups. It preserves all cycles,
boundaries and $E_2$ labels, commutes with $d_r$ for $r\ge2$, and sends
multiplication by $\lambda$ to the identity. The maps on pages are
induced by this morphism of cycle and boundary groups.

For $r\ge2$ and $r-2\le t-w$, the page map of this comparison is an
isomorphism
\[
 E_r^{s,t,w}(\nu X)\cong E_r^{s,t}(X).
\]
It preserves labels and commutes with $\lambda$ and differentials. Its
inverse gives classical lifts with the same compatibilities. In
particular, these lifts invert the localization map on the stated range.

The endpoint $t-w=r-2$ is included. Put $a=t-w$: an incoming $d_i$ has
source $\lambda$-exponent $a-i+1$, so every incoming differential with
$2\le i\le r-1$ is present when $a\ge r-2$. Outgoing cycles agree under
localization, including at the endpoint.

\lean{KIP126.Literature.Standard.inversion_data,KIP126.Literature.Standard.inversion_e2,KIP126.Literature.Standard.inversion_differential,KIP126.Literature.Standard.inversion_lambda,KIP126.Literature.Standard.inversionData,KIP126.Literature.Standard.inversionConstruction,KIP126.Literature.Standard.free_component,KIP126.Literature.Standard.free_e2,KIP126.Literature.Standard.free_lambda,KIP126.Literature.Standard.free_differential,KIP126.Literature.Standard.freeData,KIP126.Literature.Standard.classicalLift,KIP126.Literature.Standard.classicalLift_differential,KIP126.Literature.Standard.classicalLift_lambda,KIP126.Literature.Standard.free_inversion}
\end{proposition}


\begin{proposition}[Finite-page quotient comparisons and support]
\label{sem:Literature:finite-quotient-labels}
\label{sem:Literature:vanishing-and-degeneration}
\level{coarse}
\uses{sem:Literature:nu-page-coordinates,sem:AdamsSynthetic:lambda-quotients,sem:Literature:free-lambda-range}
The finite-quotient formulas follow from \cite[Theorem A.8 and Corollary A.11]{BHS}.
For every $q>0$, the finite quotient has
\[
 E_2^{s,t,w}(\nu X/\lambda^q)\cong
 \begin{cases}E_2^{s,t}(X),&0\le t-w<q,\\0,&\text{otherwise}.
 \end{cases}
\]
In this strip, quotient projection preserves labels and is an
isomorphism on $E_2$. For $r\ge2$ in the smaller range
\[
 0\le t-w,\qquad t-w+r-2<q,
\]
it induces an isomorphism
\[
 E_r^{s,t,w}(\nu X)\longrightarrow E_r^{s,t,w}(\nu X/\lambda^q).
\]
On every page $r\ge2$, one also has
\[
 E_r^{s,t,w}(\nu X)=0\quad(t<w),\qquad
 E_r^{s,t,w}(\nu X/\lambda^q)=0\quad(t-w<0\text{ or }q\le t-w).
\]

\lean{KIP126.Literature.Standard.finite_quotient_component,KIP126.Literature.Standard.finite_quotient_component_retained,KIP126.Literature.Standard.finite_quotient_retained,KIP126.Literature.Standard.finite_quotient,KIP126.Literature.Standard.finiteQuotientE2Construction,KIP126.Literature.Standard.above_diagonal,KIP126.Literature.Standard.finite_quotient_vanishing}
\end{proposition}
\begin{proof}
In the strip of BHS Corollary A.11, the quotient projection identifies
the two $E_2$ pages with the same classical labels. In the smaller
range on page $r$, the outgoing terms used to form that page remain
in the strip. Incoming terms either lie in the strip or vanish on
both sides below the diagonal. The projection is therefore still an
isomorphism there.
The support assertions follow from the zero groups on $E_2$ by passage
to successive homology pages. The calculation is entirely on the finite pages and their cycle and boundary quotients.
\end{proof}

\paragraph*{The lambda-tower convention}
For a synthetic object $Y$, repeated multiplication by $\lambda$ gives
the tower
\[
 \cdots\longrightarrow\Sigma^{0,-2}Y
 \xrightarrow{\lambda}\Sigma^{0,-1}Y\xrightarrow{\lambda}Y,
\]
whose arrows are the corresponding shifts of $\lambda_Y$ and whose
successive cofibers are shifts of $Y/\lambda$. Its exact couple gives
the lambda-Bockstein spectral sequence. In tridegree $(s,t,w)$ the
tower index is $t-w$; its initial term is
$\pi_{t-s,t}(Y/\lambda)$ for $t-w\ge0$, and zero otherwise.
We call this initial page $E_2$, so the differential convention is
$d_r=d_{r-1}^{\mathrm{tower}}$ for $r\ge2$: the tower's first
differential appears here as $d_2$. The comparison with the Adams
sequence is assumed in the following proposition.

\begin{proposition}[Adams--lambda comparison and labeled differentials]
\label{sem:Literature:adams-lambda-comparison-input}
\label{sem:Literature:differential-rigidity}
\level{coarse}
\uses{sem:AdamsSynthetic:families,sem:AdamsSynthetic:lambda-quotients,sem:AdamsSynthetic:e2-comparison,sem:Literature:free-lambda-range,sem:Literature:nu-page-coordinates}
Assume the following comparison for every spectrum $X$ in the completed
category $\mathrm{Sp}$. There are additive first-quotient equivalences
\[
 E_2^{s,t}(X)\cong\pi_{t-s,t}(\nu X/\lambda)
 \qquad(s,t\in\mathbb Z).
\]
For every integer $a$, there are morphisms, starting on page two, in both
directions between the Adams sequence of $\Sigma^{0,a}\nu X$
and the lambda-tower sequence of $\nu X$,
reindexed only in weight by $w\mapsto w-a$.
Their composites are the identity on $E_2$. On the Adams objects fixed above, require that the first-quotient equivalence agree with the
earlier label map, and that the forward morphism carry each $E_2$ label in degree
$(s,t,t+a-k)$ to that same first-quotient label, for every $k\ge0$.

Assume also the following differential comparison for every $X$, all
integer $s,t$ and $r,k\in\mathbb N$ with $r\ge2$: the classical
relation $d_r(x)=y$ is equivalent to the
lambda-tower relation between their first-quotient labels in degrees
$(s,t,t-k)$ and $(s+r,t+r-1,t-k)$, the latter using exponent $k+r-1$.
These morphisms preserve cycle and boundary groups and the
representatives on each finite page. The morphisms compare the two ambient-module presentations.

The Adams--Bockstein comparison applies to every spectrum in the completed
category, allowing zero labels; its source is
\cite[Theorem 2.5(4)]{Marek2025}; see also \cite[Theorem A.8]{BHS}
and \cite[Theorem 3.8]{LWX126}. The first-quotient label identifications
are part of the assumed comparison.

In these comparisons the classical relation $d_r(x)=y$, for $r\ge2$,
is equivalent to $d_r(x)=\lambda^{r-1}y$ in synthetic labeled notation.
More generally, the corresponding weighted relation is
\[
 d_r(\lambda^kx)=\lambda^{k+r-1}y\qquad(k\in\mathbb N).
\]
After a shift $(0,a)$, both sides have weight $t+a-k$.
These relations use the specified representatives, allowing zero
labels and earlier boundaries; see \cite[Theorem A.8(3)]{BHS}.

\lean{KIP126.Literature.Paper38ComparisonInput.adamsLambdaComparison_input,KIP126.Literature.Standard.differential_rigidity}
\end{proposition}
\begin{proof}
The page maps are inverse because they are inverse on $E_2$. Transport
the two representatives through them. The weight reindexing sends
$(s,t,t+a-k)$ to $(s,t,t-k)$ and preserves differential degree. Its
compatibility on the initial page identifies the transported labels with the
first-quotient labels of $x,y$. Apply the classical-to-lambda relation,
whose target exponent is $k+r-1$, and transport back. This gives the
synthetic relation for all stated integer shifts and all $r,k$ in range.
\end{proof}

\subsection{Infinite-page windows, shifts and restriction}
At weight $w$, only a finite range of incoming classical boundaries is
visible. For a finite $\lambda$ quotient, only a finite range of outgoing
differentials is visible as well. The following two assumptions express
these windows and specify which classical cycles represent their
classes. The convergence hypotheses in the next section interpret these
infinite-page windows in homotopy.

% AimPaper: 3.11 proposition 841 847
\begin{proposition}[The infinity page of nu X and its common labels]
\label{sem:Literature:infinite-nu-window}
\label{paper:3.11}
\label{prop:30e8b746}
\level{coarse}
\source{AimPaper Proposition 3.11}
\uses{sem:AdamsSynthetic:families,sem:SpectralSequence:cycles-boundaries,sem:SpectralSequence:permanent,sem:Literature:nu-page-coordinates,sem:Literature:vanishing-and-degeneration,paper:3.6}
Compare \cite{BHS}, Corollary A.9, and \cite{LWX126}, Notation 3.10
and Proposition 3.11 (\href{paper/main.tex.html#L841}{source text}).
Let $Z_r^{s,t}(X),B_r^{s,t}(X)\subseteq E_2^{s,t}(X)$ be the classical
cycle and boundary subgroups and put $Z_\infty=\bigcap_rZ_r$.
A permanent representative is an element of $Z_\infty$, including zero.

Assume that for every Adams object $X$ fixed above and all integers
$s,t,w$ there are additive equivalences
\[
 E_\infty^{s,t,w}(\nu X)\cong
 \begin{cases}
 Z_\infty^{s,t}(X)/B_{1+t-w}^{s,t}(X),&w\le t,\\
 0,&t<w.
 \end{cases}
\]
Assume also that this equivalence identifies representatives as follows.
For $X$ in the bounded-below finite-type range, $w\le t$,
$x\in E_2^{s,t}(X)$ and
$z\in E_\infty^{s,t,w}(\nu X)$, the weight-$w$ label of $x$ represents
$z$ if and only if $x\in Z_\infty^{s,t}(X)$ and the displayed equivalence
sends $z$ to the class of $x$ modulo $B_{1+t-w}^{s,t}(X)$.
Zero and boundary labels are allowed in this equivalence.

\lean{KIP126.Literature.Standard.eInfinity_nu_labels,KIP126.Literature.Standard.infinityNuWindow,KIP126.Literature.Standard.infinityNuLinear,KIP126.Literature.Standard.nuInfinityConstruction,KIP126.Literature.Standard.eInfinity_nu}
\end{proposition}
\begin{proof}
Use the equivalence with its assumed compatibility on representatives.
With $k=t-w\ge0$, the formula reads
$Z_\infty^{s,t}/B_{k+1}^{s,t}$: a classical permanent cycle becomes
zero precisely when it is among the boundaries visible at this weight.
The homotopy interpretation uses convergence in addition to this formula.
\end{proof}

% AimPaper: 3.12 proposition 849 855
\begin{proposition}[Finite-quotient infinity pages, labels and lambda maps]
\label{sem:Literature:infinite-finite-window}
\label{paper:3.12}
\label{prop:59f111f}
\level{coarse}
\source{AimPaper Proposition 3.12}
\uses{sem:Literature:infinite-nu-window,sem:Literature:finite-quotient-labels,sem:Literature:vanishing-and-degeneration,paper:3.6,paper:3.3}
Compare \cite{BHS}, Corollary A.11, and \cite{LWX126}, Proposition 3.12
and Remark 3.14 (\href{paper/main.tex.html#L849}{source text}).
Assume that for every Adams object $X$ fixed above, every positive
length $q$ and all integers $s,t,w$ there are additive equivalences
\[
 E_\infty^{s,t,w}(\nu X/\lambda^q)\cong
 \begin{cases}
 Z_{q-t+w}^{s,t}(X)/B_{1+t-w}^{s,t}(X),&0\le t-w<q,\\
 0,&\text{otherwise}.
 \end{cases}
\]
Require the following compatibilities with representatives and maps.
\begin{enumerate}
\item For $X$ in the bounded-below finite-type range, $0\le t-w<q$,
$x\in E_2^{s,t}(X)$ and $z\in E_\infty^{s,t,w}(\nu X/\lambda^q)$, the
projected weight-$w$ label of $x$ represents $z$ if and only if
$x\in Z_{q-t+w}^{s,t}(X)$ and the window image of $z$ is the class of
$x$ modulo $B_{1+t-w}^{s,t}(X)$.
\item For $0\le k<q$ and $0\le t-w<q-k$, the finite $\lambda^k$ map
from length $q-k$, weight $w$, to length $q$, weight $w-k$, is carried to
\[
 Z_{q-k-t+w}/B_{1+t-w}\longrightarrow
 Z_{q-t+w-k}/B_{1+t-w+k}.
\]
All groups here have bidegree $(s,t)$. The cycle indices agree and the
boundary inclusion follows from the stated ranges.
\end{enumerate}
The formula covers every $q>0$, including $q=1$, as in
\cite[Corollary A.11]{BHS}. The range $q\ge2$ is also recorded in
\cite[Proposition 3.12]{LWX126}.

\lean{KIP126.Literature.Standard.eInfinity_finite_labels,KIP126.Literature.Standard.infinityFiniteWindow,KIP126.Literature.Standard.infinityFiniteLinear,KIP126.Literature.Standard.finiteInfinityConstruction,KIP126.Literature.Standard.eInfinity_finite}
\end{proposition}
\begin{proof}
For $k>0$, the finite $\lambda^k$ map is the shift followed by the map
in the quotient triangle; for $k=0$, it is the identity. The assumed
second compatibility identifies its action on the windows with
enlargement of the boundary subgroup.
Putting $k=t-w$ rewrites
the retained quotient as $Z_{q-k}/B_{k+1}$ for $0\le k<q$.
The same description includes length one. Its interpretation in
homotopy uses the convergence and separation assumptions fixed above.
\end{proof}

\begin{proposition}[Shift, lambda and restriction on the infinity windows]
\label{sem:Literature:infinite-shift}
\label{sem:Literature:infinite-restriction}
\level{coarse}
\uses{sem:AdamsSynthetic:ess-shifts,sem:Literature:infinite-nu-window,sem:Literature:infinite-finite-window,sem:AdamsSynthetic:weight-representative-input,sem:Literature:lambda-inversion}
Compare \cite{LWX126}, Remark 3.14.
For $k\ge0$, the shift $(0,-k)$ gives filtered and $E_\infty$
comparisons natural in morphisms of the synthetic objects fixed above.
Composing with $\lambda^k$ gives weight-lowering maps, which are the
identity at $k=0$. In the classical window coordinates they act
as follows.
\begin{enumerate}
\item On $\nu X$, for $w\le t$, the map enlarges the denominator
$B_{1+t-w}$ to $B_{1+t-w+k}$.
\item On finite quotients, for $k<q$ and $0\le t-w<q-k$, it is the
natural map from the classical cycle/boundary quotient of length $q-k$,
weight $w$, to that of length $q$, weight $w-k$.
\item For $0<i\le j$ and $0\le t-w<i$, restriction
$\nu X/\lambda^j\to\nu X/\lambda^i$ is the natural map of the two
cycle/boundary quotients. Projection from $\nu X$ to a positive finite
quotient carries a permanent-cycle quotient class to its finite-cycle
quotient class in the common strip.
\end{enumerate}

\lean{KIP126.Literature.Standard.infinityShift_comparison,KIP126.Literature.Standard.infinityShift_natural,KIP126.Literature.Standard.infinityShiftData,KIP126.Literature.Standard.infinityLambda,KIP126.Literature.Standard.infinityFiniteLambda,KIP126.Literature.Standard.infinityLambda_zero,KIP126.Literature.Standard.infinityFiniteLambda_zero,KIP126.Literature.Standard.eInfinity_lambda_nu,KIP126.Literature.Standard.eInfinity_lambda_finite,KIP126.Literature.Standard.eInfinity_rho_finite,KIP126.Literature.Standard.eInfinity_rho_nu}
\end{proposition}
\begin{proof}
The filtered shift comparison and its naturality give the first maps.
Compose with powers of $\lambda$ and use the window
comparisons to obtain the stated boundary inclusions. The source and
target lengths in the finite case are $q-k$ and $q$, respectively;
the identity at zero uses shift coherence. For restriction and
projection, the label-compatible windows identify the common $E_2$
labels and their differentials under lambda inversion. The free-range
page map is an isomorphism, and the representative equations then give
the natural maps on the indicated cycle/boundary quotients.
\end{proof}


\begin{theorem}[Geometric boundary of a finite lambda lift]
\label{sem:Literature:lambda-geometric-boundary-input}
\level{coarse}
\uses{sem:AdamsSynthetic:lambda-quotients}
Assume the following description of the lambda-Bockstein boundary.
Let $X$ be any spectrum in the completed category $\mathrm{Sp}$,
$q\in\mathbb N$ with $q>0$, and
$s,t\in\mathbb Z$. For every finite lift
\[
 a\in\pi_{t-s,t}(\nu X/\lambda^q),
\]
its restriction to the first quotient and its image under the
quotient-triangle boundary
\[
 \nu X/\lambda^q\longrightarrow\Sigma^{1,-q}(\nu X/\lambda)
\]
represent a $d_{q+1}$ relation in the lambda-tower spectral sequence,
from degree $(s,t,t)$ to $(s+q+1,t+q,t)$. Both labels use the
first-quotient coordinates above, and the target uses the
suspension regrading of that boundary. The assertion applies separately to every lift, with its specified
first-quotient source and target, allowing zero and earlier boundaries.

The formula applies to every finite lift above. The defining exact couple of the lambda-Bockstein sequence computes the
next differential by this connecting morphism. Changing a finite lift
with its first-quotient class fixed changes the target by an earlier
boundary. Thus the relation holds for every such lift. See the cofiber
diagram in the proof of
\cite[Theorem 9.19(1)]{BHS} and \cite[Theorem 3.8]{LWX126}.
\lean{KIP126.Literature.Paper38BoundaryInput.lambda_geometric_boundary_input}
\end{theorem}

\subsection{Finite lifts and the Bockstein boundary}
The next result connects continuation of a classical label with a lift to
a finite quotient. Its boundary conclusion gives a suitable lift realizing the represented
page differential.

% AimPaper: 3.8 theorem 813 823
\begin{theorem}[Finite lifts and the Bockstein differential]
\label{sem:Literature:finite-bockstein}
\source{AimPaper Theorem 3.8}
\level{coarse}
\uses{sem:AdamsSynthetic:lambda-quotients,sem:AdamsSynthetic:e2-comparison,sem:Literature:finite-quotient-labels,sem:Literature:adams-lambda-comparison-input,sem:Literature:lambda-geometric-boundary-input,sem:Literature:infinite-finite-window}
The finite lifting criterion is \cite[Theorem 3.8]{LWX126}; see also
\cite[Theorem A.1(1) and Corollary A.11]{BHS}.\par

Let $X$ be in the bounded-below finite-type spectrum range and $s,t\in\mathbb Z$. Use the first-quotient identification
\[
 E_2^{s,t}(X)\cong\pi_{t-s,t}(\nu X/\lambda).
\]
For each positive integer $q$ and $x\in E_2^{s,t}(X)$, write $\bar x$ for its first-quotient image and $\rho_q$ for restriction to the first quotient. Then
\[
 x\text{ reaches }E_{q+1}
 \quad\Longleftrightarrow\quad
 \exists a\in\pi_{t-s,t}(\nu X/\lambda^q),\quad\rho_q(a)=\bar x.
\]
When these conditions hold, such an $a$ can be chosen so that its Bockstein label represents $d_{q+1}(x)$ in bidegree $(s+q+1,t+q)$. The label is obtained from the quotient-triangle boundary
\[
 \nu X/\lambda^q\longrightarrow\Sigma^{1,-q}(\nu X/\lambda)
\]
and the same first-quotient identification. This is the relation on $E_{q+1}$ represented by its specified $E_2$
boundary label. The lift criterion allows zero labels and earlier boundaries.

\lean{KIP126.Literature.Standard.finite_lift_criterion,KIP126.Literature.Standard.bockstein_differential}
\end{theorem}
\begin{proof}
\uses{sem:Literature:finite-quotient-labels,sem:AdamsSynthetic:selected-convergence}
At weight $t$, the finite-quotient infinity-page description identifies
its lowest relevant filtration piece with the classical cycles that
reach $E_{q+1}$. The finite filtration window, convergence and separation
provide a homotopy representative of each such infinity class.
Naturality of restriction to the first quotient identifies its image
with $\bar x$. Conversely, the infinity
label of a homotopy lift restricts to that same label and hence gives
the required classical cycle. This proves the lift equivalence using
restriction, with zero and earlier boundaries allowed.

For the differential assertion, the boundary assumption identifies the
restriction of a finite lift and its quotient-triangle boundary with
the two sides of a lambda-Bockstein differential. The Adams--lambda
comparison transports this relation to $d_{q+1}(x)$. Both comparisons
use the first-quotient coordinates above, which
identify the boundary label. This yields the stated lift and the
represented differential on $E_{q+1}$.
\end{proof}

\section{Realization, completeness and untruncated lifts}
Finite quotients detect how long an Adams class survives. Convergence
and separation supply untruncated synthetic lifts of permanent cycles;
a map of Adams towers identifies their realization with classical
homotopy representatives. The following results specify these comparisons.

\subsection{Maps of towers and detection under realization}

\begin{definition}[Lambda-inversion comparisons of Adams towers]
\label{sem:Literature:lambda-inversion-object-comparison}
\label{sem:Literature:lambda-inversion-e2-map}
\label{sem:Literature:lambda-inversion-adams-comparison}
\level{coarse}
\uses{sem:Literature:realization-coordinates,sem:AdamsSynthetic:selected-convergence,sem:SpectralSequence:filtration,sem:AdamsSynthetic:e2-comparison,sem:AdamsSynthetic:normalized-map}
Compare \cite{Pst}, Proposition 4.40, and \cite{BHS}, Theorem A.8.
Use the realization coordinates of
Definition~\ref{sem:Literature:realization-coordinates}.

\emph{Comparison of towers.} Let $U$ be synthetic and $X$ classical.
A lambda-inversion comparison consists of isomorphisms
\[
 \theta:\operatorname{Re}(U)\xrightarrow{\sim}X,\qquad
 \theta_w:\operatorname{Re}(\Sigma^{0,-w}U)\xrightarrow{\sim}X
 \quad(w\in\mathbb Z),
\]
and maps from the realized stages and cofiber layers of the Adams tower
of $\Sigma^{0,-w}U$ to the classical tower of $X$. These commute with
all tower transitions, projections and cofiber inclusions. For every
$\beta\in\pi_{m,w}U$ they satisfy
\[
 \theta_w\circ\operatorname{Re}(\mathrm{wt}_{m,w}(\beta))\circ r_m
 =\theta\circ\operatorname{Re}(\beta)\circ\rho_{m,w},
\]
where
\[
 \mathrm{wt}_{m,w}:\pi_{m,w}U\xrightarrow{\sim}
 [S^m_{\mathrm{syn}},\Sigma^{0,-w}U]
\]
is weight desuspension. Require also that the
induced maps of cycle and boundary groups preserve every finite and
infinite level, with the first-page map induced by realization,
$r_{t-s}$ and the cofiber-layer map.

\emph{Page-two maps.} For such a comparison $K$, write
\[
 K_{2,w}^{s,t}:E_2^{s,t,w}(U)\longrightarrow E_2^{s,t}(X)
\]
for the induced map on page two.

\emph{The synthetic analogues.} Assume that these comparisons exist as
$K_X$ for $\nu X$ and $K_{X,a}$ for $\Sigma^{0,a}\nu X$, for every
Adams object $X$ fixed above and every $a\in\mathbb N$.
The target isomorphism of $K_X$ is the recovery isomorphism
$\theta_X:\operatorname{Re}(\nu X)\cong X$. Its shifted counterpart
satisfies
\[
 \theta_{X,a}\circ\operatorname{Re}(\lambda^{+a}_{\nu X})=\theta_X,
 \qquad \lambda^{+a}_{\nu X}:\nu X\longrightarrow\Sigma^{0,a}\nu X.
\]
Here $\lambda^{+a}_U$ is the shifted map
$\Sigma^{0,a}(\lambda_U^a)$, whose source is identified with $U$ by
shift coherence. The plus sign records the positive target weight.
For every $s,t\in\mathbb Z$ and $k\in\mathbb N$, the maps
$(K_X)_{2,t-k}^{s,t}$ and $(K_{X,a})_{2,t+a-k}^{s,t}$ send the
unshifted and shifted synthetic images of an $E_2^{s,t}(X)$ label back to that
same label.

The assumed tower comparisons are induced by lambda inversion, with
the stage, cofiber and label compatibilities just specified. They are
comparisons in the completed categories fixed above.

\lean{KIP126.Literature.LambdaInversionAdamsObjectComparison,KIP126.Literature.LambdaInversionAdamsObjectComparison.e2Map,KIP126.Literature.LambdaInversionAdamsComparison,KIP126.Literature.Standard.lambdaInversionAdamsComparison}
\end{definition}

\begin{theorem}[Realization preserves detection along a tower comparison]
\label{sem:Literature:lambda-inversion-detection}
\level{coarse}
\uses{sem:Literature:lambda-inversion-object-comparison,sem:Literature:lambda-inversion-e2-map,sem:Literature:realization-coordinates,sem:AdamsSynthetic:selected-convergence}
Let $K$ be a tower comparison between $U$ and $X$. For any
$s,t,w\in\mathbb Z$, suppose that
$x\in E_2^{s,t,w}(U)$ detects a class
$\beta\in\pi_{t-s,w}U$. Then
$K_{2,w}^{s,t}(x)$ classically detects
\[
\theta\circ\operatorname{Re}(\beta)\circ\rho_{t-s,w}
\in\pi_{t-s}X.
\]
The assertion includes zero and boundary labels and concerns the
realization of the given class $\beta$.

\lean{KIP126.Literature.LambdaInversionAdamsObjectComparison.detects}
\end{theorem}
\begin{proof}
\uses{sem:Literature:lambda-inversion-object-comparison,sem:AdamsSynthetic:selected-convergence}
Lift $\beta$ to its synthetic Adams stage. The canonical
convergence formula identifies the cofiber image of this lift with a
common infinite-cycle representative of its detecting class. Realize
the lift and compose with the stage map. The projection and
sphere-coordinate equations give exactly
$\theta\operatorname{Re}(\beta)\rho_{t-s,w}$, while the cofiber square
and the ambient comparison identify its layer class with the image of
the original infinite representative. Preservation of cycles and
boundaries identifies its page-two label with $K_{2,w}^{s,t}(x)$.
The classical canonical convergence formula now gives the claimed
detection.
\end{proof}

\subsection{Complete sources and synthetic lifts}
\label{sem:Literature:lambda-completeness}

\paragraph*{Nilpotent completeness and lambda completeness}
The inverse-limit observation in \cite{LWX126}, Notation 4.1, refers to
\cite{BHS}, Definition A.12 and Proposition A.13, specialized to
$E=H\mathbb F_2$. In that source category, the following conditions
on a spectrum $X$ are equivalent: $X$ is $H\mathbb F_2$-nilpotent
complete; $\nu X$ is $\nu H\mathbb F_2$-nilpotent complete; and
$\nu X$ is $\lambda$-complete. The last condition means that the
canonical map
\[
\nu X\longrightarrow\operatorname*{holim}_{m\geq1}
       \bigl((S^{0,0}/\lambda^m)\otimes\nu X\bigr)
\]
is an equivalence. Under the source tensor/cofiber identification, the
objects on the right are $\nu X/\lambda^m$. The limit is taken with
the quotient restriction maps and the cone from $\nu X$.

The source proof compares the canonical classical and synthetic Adams
towers using the spectral Yoneda functor and its connective covers.
Right completeness of the homological t-structure identifies the
vanishing of their residual homotopy limits. Finite lambda quotients
are synthetically nilpotent complete, and
$\nu H\mathbb F_2\otimes\nu X$ is lambda-complete, by
\cite{BHS}, Lemmas A.14--A.15. Applying these facts to the canonical
Adams resolution gives the equivalence with lambda completeness.

For the ordinary category fixed above, we assume that nilpotent
completeness and this homotopy-limit cone agree with the constructions
in \cite{BHS}. These assumptions are additional to the finite-quotient
comparisons.

\begin{definition}[Completeness hypotheses for an Adams tower]
\label{sem:Literature:bhs-applicability}
\level{coarse}
See \cite[Definition A.12]{BHS}.\par

Write
$\cdots\xrightarrow{d_1}R_1X\xrightarrow{d_0}R_0X\to X$
for the residual classical Adams tower; the images of its stages in
homotopy define the Adams filtration $F^s\pi_nX$. Require countable
products and invertibility of the Milnor map
\[
 1-\sigma:\prod_{m\ge0}R_mX\longrightarrow\prod_{m\ge0}R_mX,
 \qquad p_m\sigma=d_m p_{m+1},
\]
where $p_m$ is the $m$th projection. This is the derived-limit condition
on the residual tower. Require also completeness of the
homotopy filtration: for every $n\in\mathbb Z$, the canonical maps
to its quotients exhibit
\[
 \pi_nX\xrightarrow{\sim}
 \varprojlim_{s\in\mathbb Z}\pi_nX/F^s\pi_nX.
\]
Finally require $\bigcap_sF^s\pi_nX=0$ in every degree, together with
the canonical convergence identifications
$\operatorname{gr}_F^s\pi_nX\cong E_\infty^{s,n+s}(X)$.
These are the complete-source hypotheses.

\lean{KIP126.Literature.Standard.AdamsComparisonApplicability}
\end{definition}

\begin{theorem}[Permanent lifts and filtration on bounded-below finite-type spectra]
\label{sem:Literature:complete-permanent-lifts}
\level{coarse}
\uses{sem:AdamsSynthetic:paper-spectrum-scope,sem:AdamsSynthetic:e2-comparison,sem:AdamsSynthetic:selected-convergence}
Compare \cite{BHS}, Theorem A.1 and Corollary A.19.\par

Let $X$ be in the bounded-below finite-type spectrum range. For
$s,t\in\mathbb Z$, a label $x\in E_2^{s,t}(X)$ has a common
infinite-cycle representative if and only if its first-quotient class
lifts to $\pi_{t-s,t}(\nu X)$. This includes zero labels and uses the
zero-shift first-quotient comparison. Moreover, for
$a\in\pi_{m,w}(\nu X)$ and $w-m\le s$,
\[
 a\in F^s\pi_{m,w}(\nu X)
 \quad\Longleftrightarrow\quad
 a=\lambda^{m+s-w}b
 \text{ for some }b\in\pi_{m,m+s}(\nu X).
\]

\lean{KIP126.Literature.Standard.permanent_lift,KIP126.Literature.Standard.filtration_lambda}
\end{theorem}
\begin{proof}
\uses{sem:Literature:infinite-nu-window,sem:Literature:infinite-finite-window,sem:AdamsSynthetic:finite-quotient-separation,sem:AdamsSynthetic:lambda-quotients,sem:AdamsSynthetic:selected-convergence,sem:AdamsSynthetic:e2-comparison}
First consider a permanent label $x\in E_2^{s,t}(X)$. The infinite-page
window at weight $t$ gives an infinite class on $\nu X$ with this
label. Lift its associated graded class to a class $a$ in filtration
$s$. In $\nu X/\lambda$, the infinite page vanishes above filtration
$s$ in this stem and weight. Separation makes filtration $s+1$ zero,
so naturality and first-quotient detection identify the projection
of $a$ with the first-quotient class of $x$.

Conversely, an arbitrary class $a\in\pi_{t-s,t}(\nu X)$ has
filtration at least $s$ by the above-diagonal vanishing of its infinite
page. Choose a permanent label for its associated graded class via
the same infinite window. Projection and the preceding uniqueness in
the first quotient identify this label with the given $x$, since the
first-quotient comparison is injective. This proves the
permanent-lift equivalence, including zero labels.

Put $k=m+s-w\ge0$. The infinite-page window of $\nu X$ vanishes when
$t<w$, so every class of bidegree $(m,m+s)$ has filtration at least
$s$. Precomposition by $\lambda^k$ preserves this filtration, proving
one implication.

Conversely, project a class of filtration at least $s$ to
$\nu X/\lambda^k$. If $k>0$, its infinite-page window vanishes in all
filtrations $j\ge s$, since $m+j-w\ge k$. The convergence comparison
therefore identifies all these filtration pieces. Separation of the
finite quotient makes the projected class zero. Cofiber exactness and
the fully faithful weight shift then supply a class of bidegree
$(m,m+s)$ whose $\lambda^k$ multiple is the original class. For $k=0$,
the quotient projection is zero because $\lambda^0$ is an isomorphism,
and the same exactness argument applies. Nonpositive filtrations are
covered by the zeroth tower stage and monotonicity.

The argument uses the windows, first-quotient detection, and
finite-quotient separation.
\end{proof}

\begin{theorem}[The realization kernel and lambda torsion]
\label{sem:Literature:realization-kernel}
\level{coarse}
\uses{sem:AdamsSynthetic:families,sem:AdamsSynthetic:recovery,sem:Literature:realization-coordinates}
The underlying-spectrum comparison and its homotopy colimit formula
are given in \cite[Proposition 4.40, Remark 4.41 and
Propositions 5.6--5.7]{Pst}.\par

Assume the following kernel identity for realization.
Let $Y\in\mathcal S$, where $\mathcal S$ contains the synthetic sphere and
$\nu X$ for every $X\in\mathcal C$ and is closed under bidegree shifts
and finite lambda quotients. For every $m,w\in\mathbb Z$ and
$a\in\pi_{m,w}Y$, the map $\operatorname{Re}(a)$ is zero if and only if
some finite power of $\lambda$ annihilates $a$:
\[
\operatorname{Re}(a)=0\quad\Longleftrightarrow\quad
\exists k\in\mathbb N,\ \lambda^ka=0.
\]


\lean{KIP126.Literature.Standard.selected_realization_kernel}
\end{theorem}
\begin{proof}
For the underlying-spectrum functor, \cite[Remark 4.41]{Pst} identifies
homotopy with the filtered colimit along multiplication by $\lambda$.
A class is zero in this colimit if and only if its image is zero at
some finite stage, equivalently if a finite power of $\lambda$ kills it.
Under the assumed kernel comparison of realization with this
underlying-spectrum functor, the criterion has the stated form on
$\mathcal S$.
\end{proof}

\begin{theorem}[Lifetime of a lift on bounded-below finite-type spectra]
\label{sem:Literature:complete-lifetime}
\level{coarse}
\uses{sem:AdamsSynthetic:paper-spectrum-scope,sem:AdamsSynthetic:e2-comparison}
Compare \cite{BHS}, Theorem A.1(2a).\par

Let $X$ be in the bounded-below finite-type spectrum range. Let
$r\ge1$, $x\in E_2^{s,t}(X)$, and
$a\in\pi_{t-s,t}(\nu X)$ have first-quotient label $x$ under the
zero-shift comparison. If $x$ has a common permanent representative
and survives nontrivially to $E_{r+1}$, then $\lambda^{r-1}a\ne0$.

\lean{KIP126.Literature.Standard.realization_lifetime}
\end{theorem}
\begin{proof}
\uses{sem:Literature:infinite-nu-window,sem:Literature:infinite-finite-window,sem:Literature:infinite-restriction,sem:Literature:infinite-shift,sem:AdamsSynthetic:e2-comparison,sem:AdamsSynthetic:selected-convergence,sem:AdamsSynthetic:ess-shifts}
The vanishing of the infinite page above the diagonal
places $a$ in filtration $s$. At weight $t$, projection
$\nu X\to\nu X/\lambda$ is injective on the infinite page: in the
window coordinates it is the inclusion of permanent cycles into the
finite-cycle quotient with the same boundary subgroup. The
first-quotient detection comparison and naturality of the convergence
identification therefore show that $a$ is detected by the infinite
class of the given label $x$.

Set $k=r-1$. Filtered shift and multiplication by $\lambda^k$ transport
this detection to weight $t-k$. The window comparison sends its
infinite class to
\[
[x]\in Z_\infty^{s,t}(X)/B_{1+t-(t-k)}^{s,t}(X)
     =Z_\infty^{s,t}(X)/B_r^{s,t}(X).
\]
If $\lambda^k a=0$, its associated graded class is zero, so
$x\in B_r^{s,t}(X)$. Such a label represents zero on $E_{r+1}$,
contradicting the supplied nonzero continuation. When $r=1$,
$\lambda^0$ is the identity; the same argument
at weight $t$ contradicts nonzero survival on $E_2$.

Thus the page on which a label becomes a boundary bounds the lambda
torsion of its lift. The argument uses the windows and convergence.
\end{proof}

\begin{theorem}[Detection and prescribed realization on bounded-below finite-type spectra]
\label{sem:Literature:complete-detection}
\level{coarse}
\uses{sem:AdamsSynthetic:paper-spectrum-scope,sem:Literature:realization-coordinates,sem:AdamsSynthetic:e2-comparison,sem:AdamsSynthetic:selected-convergence}
These are the detection and prescribed-realization statements of \cite[Theorem A.1(2b),(3b)]{BHS}.\par

Let $X$ be in the bounded-below finite-type spectrum range, with the
Adams and synthetic comparisons above. Let $x\in E_2^{s,t}(X)$ have
a common infinite-cycle representative with nonzero image in
$E_\infty$. Assume the following two realization properties:
\begin{enumerate}
\item Every lift in $\pi_{t-s,t}(\nu X)$ with first-quotient label
$x$ realizes to a classical class detected by $x$.
\item For every classical class $a\in\pi_{t-s}X$ detected by $x$,
there is a lift in $\pi_{t-s,t}(\nu X)$ with first-quotient label $x$
and realization $a$.
\end{enumerate}
Both assertions use the zero-shift and realization comparisons.


\lean{KIP126.Literature.Standard.realization_detection,KIP126.Literature.Standard.realization_prescribed_lift}
\end{theorem}
\begin{proof}
Apply \cite[Theorem A.1(2b),(3b)]{BHS} through the realization and
first-quotient comparisons fixed above, with the completeness and
convergence hypotheses fixed above. The first property concerns every
synthetic lift, and the second provides a lift of every classical class
with the stated detector.
\end{proof}

\begin{theorem}[A boundary has a suitably torsion lift]
\label{sem:Literature:complete-boundary-lift}
\level{coarse}
\uses{sem:AdamsSynthetic:paper-spectrum-scope,sem:AdamsSynthetic:e2-comparison}
Compare \cite{BHS}, Theorem A.1(3a).\par

Let $X$ be in the bounded-below finite-type spectrum range. If $r\ge2$
and $y\in E_2^{s,t}(X)$ has a common infinite-cycle representative
and is hit on page $r$, then there is a lift
$a\in\pi_{t-s,t}(\nu X)$ with first-quotient label $y$ under the
zero-shift comparison and $\lambda^{r-1}a=0$.

\lean{KIP126.Literature.Standard.realization_boundary_lift}
\end{theorem}
\begin{proof}
\uses{sem:Literature:complete-permanent-lifts,sem:Literature:infinite-nu-window,sem:Literature:infinite-shift,sem:AdamsSynthetic:e2-comparison,sem:AdamsSynthetic:selected-convergence,sem:AdamsSynthetic:ess-shifts,sem:AdamsSynthetic:lambda-quotients}
Put $k=r-1>0$ and choose a lift
$\alpha\in\pi_{t-s,t}(\nu X)$ of the given permanent label $y$.
The differential hitting $y$ places it in $B_r^{s,t}(X)$.
Consequently $\lambda^k\alpha$ has zero infinite-page class in
weight $t-k$, and detection raises its filtration to at least $s+1$.
The filtration--$\lambda$ comparison now supplies
$\beta\in\pi_{t-s,t+1}(\nu X)$ with
\[
\lambda^{k+1}\beta=\lambda^k\alpha.
\]
Set $a=\alpha-\lambda\beta$. Cofiber exactness makes the
first-quotient image of $\lambda\beta$ zero, so $a$ has the same
label $y$. Coherence and additivity of multiplication by $\lambda$
give $\lambda^k a=0$.

Correcting a lift by a lambda multiple gives the required torsion.
\end{proof}


\section{Cofibers, filtered lifts and normalized triangles}
We next compare the arrows of classical and synthetic cofiber triangles.
Pstr\k{a}gowski's cofiber criterion and the filtered-lift theorem of BHS
explain when a classical triangle lifts synthetically. Normalization
records the filtration in its weight shifts. Later finite-quotient
arguments need the resulting arrows and their boundary signs.

\subsection{Preserved cofiber triangles and full lifts}

% AimPaper: 3.1 proposition 759 767
\begin{proposition}[Synthetic cofiber sequences and homology exactness]
\label{paper:3.1}
\source{AimPaper Proposition 3.1}
\level{coarse}
This is \cite[Lemma 4.23]{Pst}; see also
\cite[Proposition 3.1]{LWX126}
(\href{paper/main.tex.html#L759}{original statement}).\par
\label{prop:1f7950df}
    Suppose that
    $$X\fto{f} Y\fto{g} Z$$
    is a cofiber sequence of spectra. Then
    $$\nu X\fto{\nu f}\nu Y\fto{\nu g} \nu Z$$
    is a cofiber sequence of ${\HF}$-synthetic spectra if and only if
    $$0\to {\HF}_*X\fto{{\HF}_*f} {\HF}_*Y\fto{{\HF}_*g} {\HF}_*Z\to 0$$
    is a short exact sequence.

\uses{sem:AdamsSynthetic:nu,sem:AdamsSynthetic:homology-exact}


\lean{KIP126.Literature.Standard.nu_cofiber}

\end{proposition}
\begin{proof}
Specialize \cite{Pst}, Lemma 4.23, to $E=H\mathbb F_2$. It identifies
the cofiber sequences preserved by the synthetic analogue functor
precisely with those whose $E$-homology sequence is short exact. Apply
this to $f$ and $g$. The statement requires no connectivity or
finiteness assumption.
\end{proof}

% AimPaper: 3.16 proposition 912 919
\begin{proposition}[Synthetic lifts and Adams filtration]
\label{paper:3.16}
\source{AimPaper Proposition 3.16}
\level{coarse}
The factorization is \cite[Lemma 9.15]{BHS}; see
\cite[Proposition 3.16]{LWX126}
(\href{paper/main.tex.html#L912}{original statement}).\par
\label{prop:ef21f9bc}
    If a map $f:X\to Y$ has Adams filtration $\AF(f)=k$, then there exists a factorization
    $$\xymatrix{
     & \Sus^{0,-k}\nu Y\ar[d]^{\lambda^k}\\
     \nu X \ar[r]_{\nu f}\ar@{-->}[ru]^{\Sus^{0,-k}\tilde f} & \nu Y
    }$$
    where $\tilde f: \Sus^{0,k}\nu X\to \nu Y$ is called a synthetic lift of $f$.

\uses{paper:3.3,sem:AdamsSynthetic:filtration-comparison}


\lean{KIP126.Literature.Standard.lift}

\end{proposition}
\begin{proof}
Apply \cite{BHS}, Lemma 9.15, to the given map of Adams filtration $k$. Its filtration-$k$ lift to the Adams tower corresponds in the synthetic category to a factorization of $\nu f$ through multiplication by $\lambda^k$. Translating the shift to the present weight convention gives $\widetilde f:\Sigma^{0,k}\nu X\to\nu Y$ and the required equality $\nu f=\lambda^k\widetilde f$, with suspensions inserted as displayed. The lift is an existence statement, with the usual ambiguity in
division by $\lambda^k$.
\end{proof}

% AimPaper: 3.17 proposition 929 937
\begin{proposition}[Synthetic lifts of distinguished triangles]
\label{paper:3.17}
\source{AimPaper Proposition 3.17}
\level{coarse}
See \cite[Proposition 3.17]{LWX126}
(\href{paper/main.tex.html#L929}{original statement}) and the proof
of \cite[Lemma 9.15]{BHS}.\par

\label{prop:41561db2}
    Suppose that we have a  distinguished triangle of spectra
    $$X\fto{f} Y\fto{g} Z\fto{h} \Sus X$$
    with $\AF(h)>0$, and consequently a short exact sequence on $\HF$-homology
    $$0\to H_*X\fto{H_*f} H_*Y\fto{H_*g} H_*Z\to 0.$$
    Then there exists a distinguished triangle of synthetic spectra
    $$\nu X\fto{\nu f}\nu Y\fto{\nu g} \nu Z\fto{\Sus^{0,-1}\hat h} \Sus^{0,-1}\nu \Sus X=\Sus^{1,0}\nu X$$
    such that $\nu h=\lambda \hat h$.

\uses{paper:3.1,paper:3.3}


\lean{KIP126.Literature.Standard.triangle_lift}

\end{proposition}
\begin{proof}
The construction in the proof of \cite[Lemma 9.15]{BHS} gives the
connecting map with both the required factorization and the
distinguished-triangle property; see also the
\href{paper/main.tex.html#L938}{proof of LWX Proposition 3.17}.
\end{proof}


\subsection{Weight shifts of normalized arrows}
The normalized-triangle theorem places the filtration exponents in
the weight shifts. The Hopf-$\nu$ cofiber is the concrete triangle used
later.

% AimPaper: 3.20 proposition 965 971
\begin{proposition}[Distinguished triangles with filtration shifts]
\label{paper:3.20}
\source{AimPaper Proposition 3.20}
\level{coarse}
See \cite[Proposition 3.20]{LWX126}
(\href{paper/main.tex.html#L965}{original statement}).\par

Suppose that we have a distinguished triangle of spectra
    $$X\fto{f} Y\fto{g} Z\fto{h} \Sus X$$
    such that $e(f)+e(g)+e(h)=1$. Then we have a distinguished triangle of synthetic spectra
    $$\nu X\fto{\hat f} \Sus^{0,-e(f)}\nu Y\fto{\hat g} \Sus^{0,-e(f)-e(g)}\nu Z\fto{\hat h} \Sus^{0,-1}\nu \Sus X$$
    (We suppress shifts in the arrow names while displaying them on the objects.)

\uses{paper:3.17,sem:AdamsSynthetic:normalized-map}


For the bounded-below finite-type spectrum range, the same three
normalized arrows give the equivalent positive-source form
\[
\Sigma^{0,e(f)}\nu X\xrightarrow{F}\nu Y
\xrightarrow{G}\Sigma^{0,-e(g)}\nu Z
\xrightarrow{H}\Sigma^{1,e(f)}\nu X.
\]
This is the displayed triangle shifted by $\Sigma^{0,e(f)}$, using
$\nu\Sigma X\cong\Sigma^{1,1}\nu X$ at the last vertex.
Here $F$ is the positive-source lift of $f$, $G$ is the
negative-target form of the lift of $g$, and $H$ is the
lift of $h$ transported through the suspension and weight
isomorphisms used in the first display.


\lean{KIP126.Literature.Standard.normalized_triangle_distinguished}

\end{proposition}
\begin{proof}
Exactly one of $e(f),e(g),e(h)$ equals one. Rotate the given classical triangle until its connecting map has positive Adams filtration, apply Proposition~\ref{paper:3.17}, and rotate the resulting synthetic distinguished triangle back. The identities $\nu u=\lambda^{e(u)}\widehat u$ give the accumulated weight shifts $0$, $-e(f)$, $-e(f)-e(g)$, and $-1$ at the four displayed vertices.

\end{proof}

\begin{proposition}[The normalized cofiber triangle of nu]
\label{sem:Literature:nu-cofiber}
\level{coarse}
\uses{sem:Literature:hopf-classes,sem:AdamsSynthetic:normalized-map}
This is the normalized triangle used in \cite[Example 6.19]{LWX126}.\par

For the classical triangle $S^3\xrightarrow{\nu}S^0\to C\nu\to S^4$, the normalized exponent of $\nu$ is one and those of the bottom- and top-cell maps are zero. The normalized $\nu$ has first-quotient label $h_2$, and the triangle of normalized arrows is distinguished.


\lean{KIP126.Literature.Standard.nuCofiber_nu_exponent,KIP126.Literature.Standard.nuCofiber_bottom_exponent,KIP126.Literature.Standard.nuCofiber_top_exponent,KIP126.Literature.Standard.nuCofiber_normalized_label,KIP126.Literature.Standard.nuCofiber_triangle,KIP126.Literature.Standard.nuCofiber_exponent_sum,KIP126.Literature.Standard.nuCofiber_distinguished,KIP126.Literature.Standard.normalizedNu}
\end{proposition}
\begin{proof}[Proof of the distinguished-triangle clause]
\uses{paper:3.1,paper:3.16,paper:3.17,sem:AdamsSynthetic:normalized-map,sem:Literature:hopf-classes,sem:Literature:realization-kernel,sem:AdamsSynthetic:filtration-comparison,sem:AdamsSynthetic:categories}
Rotate the cofiber triangle of the Hopf map. Its positive Adams
filtration makes the induced map on mod-two homology zero. The homology
long exact sequence therefore gives a short exact sequence on the
rotated triangle, and the rotated connecting map has filtration at
least one. The synthetic triangle-lift theorem supplies one
connecting arrow with both the required factorization and
distinguishedness. Inverse rotation gives a candidate first arrow on
the three objects.

To identify this arrow, use the suspension square for the same Hopf
map and the comparisons between the synthetic and classical
spheres. The square follows by normalizing the rotated lifts and
comparing their realizations. In the needed sphere bidegree,
realization is injective: its kernel is lambda torsion, while the
first-quotient vanishing makes the relevant lambda maps injective.
The absence of maps from the sphere back to its positive suspension
removes the remaining normalization ambiguity. The resulting arrow
agrees with the normalized Hopf map; the exponent-zero bottom
and top arrows agree under the suspension comparisons. Transporting
the inverse-rotated triangle along these identifications proves the
claim, with the rotation signs retained.

The triangle lift, realization kernel and filtration comparisons give
distinguishedness. The assumed first-quotient label is $h_2$.
\end{proof}


\section{Homology exactness and the second-page operations}
The preceding triangle comparisons lead to two kinds of second-page
exact sequence. A homology short exact triangle gives an ordinary
connecting map of degree $(1,0)$. A compatible normalized triangle gives
a first-quotient operation of degree $(1,1)$ after the weight shifts.
We distinguish these operations by their constructions and domains.
The degree-$(1,1)$ operation is defined through the first-quotient triangle.

\subsection{Homology short exact triangles}

\begin{theorem}[The ordinary $E_2$ sequence of a homology short exact triangle]
\label{sem:Literature:spectrum-homology-short-exact-e2}
\level{coarse}
\source{Derived from the ordinary coefficient Adams tower}
\uses{sem:AdamsSynthetic:homology-exact,sem:AdamsSynthetic:families}
Let
\[
  A\xrightarrow{f}B\xrightarrow{g}C\xrightarrow{h}\Sigma A
\]
be a distinguished triangle in the $H\mathbb F_2$-local classical
category $\mathrm{Sp}$. Suppose that, for every
$n\in\mathbb Z$, its first two maps give a short exact sequence
\[
  0\longrightarrow H_n(A;\mathbb F_2)
  \xrightarrow{H_n(f)}H_n(B;\mathbb F_2)
  \xrightarrow{H_n(g)}H_n(C;\mathbb F_2)\longrightarrow0.
\]
There is an additive connecting map
\[
  \partial^{s,t}:E_2^{s,t}(C)\longrightarrow E_2^{s+1,t}(A)
\]
for which the following sequence is exact, for all $s,t\in\mathbb Z$:
\[
  \cdots\longrightarrow E_2^{s,t}(A)
  \xrightarrow{f_2}E_2^{s,t}(B)
  \xrightarrow{g_2}E_2^{s,t}(C)
  \xrightarrow{\partial^{s,t}}E_2^{s+1,t}(A)
  \xrightarrow{f_2}E_2^{s+1,t}(B)\longrightarrow\cdots .
\]
Here $f_2$ and $g_2$ are the maps induced on the ordinary Adams page.
The proof constructs $\partial^{s,t}$ as the boundary in cohomology.

\lean{KIP126.Literature.Standard.E2Exactness.spectrumHomologyShortExactConnectingHom,KIP126.Literature.Standard.E2Exactness.spectrumHomologyShortExactLongExactSequence}
\end{theorem}
\begin{proof}
\uses{sem:Literature:connected-homology,sem:AdamsSynthetic:completed-coproducts,sem:AdamsSynthetic:families,sem:AdamsSynthetic:coefficient-tower-e2-cohomology}
Let $L_s$ be the successive layers of a coefficient Adams tower of the
sphere, and put $K_s=\Sigma^sL_s$. The connecting maps define a
cochain differential $K_s\to K_{s+1}$. Tensoring with $X$ and taking
$\pi_t$ gives a functorial cochain complex $C_t(X)$ with
$H^s(C_t(X))\cong E_2^{s,t}(X)$.

We check that the homology short exact triangle gives a degreewise
short exact sequence of these complexes. The coefficient unit splits
after tensoring with $H\mathbb F_2$, by the universal-coefficient
comparison and the sphere-coproduct construction. Thus maps into a
coefficient layer extend across coefficient-null maps. Degreewise
homotopy bases, preservation of coproducts and Whitehead's theorem
also show that a map induced by tensoring with the coefficient
spectrum is zero when it vanishes on every homotopy group.

The hypothesis makes $H\mathbb F_2\otimes h$ zero on homotopy, hence
zero. Tensoring with every shifted Adams layer therefore kills $h$;
the negative layers are zero. It follows that
$C_t(A)\to C_t(B)\to C_t(C)$ is degreewise short exact. Its cohomology
long exact sequence gives $\partial^{s,t}$.

Finally, these cohomology groups identify naturally with the stated
Adams pages. To check the maps, compare the shifted-layer cochain maps
at their degree-zero augmentation. Exactness of the cofiber triangles
and extension into coefficient layers supply a cochain homotopy, so
the induced maps on cohomology are $f_2$ and $g_2$ in every integer
degree. Transporting the cohomology long exact sequence proves the claim.
\end{proof}

\subsection{Normalized finite quotient rows and their exact sequence}
For the normalized operations we first need a distinguished row with
the quotient arrows and the correct boundary sign. Its
quotient-length-one instance then gives the displayed exact sequence.

\begin{definition}[The finite quotient comparison and its exact rows]
\label{sem:Literature:finite-quotient-comparison}
\label{sem:Literature:finite-quotient-horizontal-row}
\level{coarse}
\uses{sem:AdamsSynthetic:categories,sem:AdamsSynthetic:lambda-quotients,paper:3.3}
Compare \cite{LWX126}, Definition 3.3 and the smash construction in the
proof of Theorem 6.12.\par

For each $q>0$, let
\[
Q_qX=\operatorname{cofib}(\lambda_X^q),\qquad
C_q=Q_qS^{0,0},\qquad \iota_X:X\longrightarrow Q_qX.
\]
Let $Q_qf$ be the induced map on these cofibers. Assume coherent
suspension isomorphisms
\[
\varepsilon_{n,X}:Q_q(\Sigma^nX)\xrightarrow{\sim}\Sigma^nQ_qX
\qquad(n\in\mathbb Z)
\]
making $Q_q$ an exact functor, and a natural isomorphism
\[
\beta_X:Q_qX\xrightarrow{\sim}C_q\otimes X.
\]
Thus, for every $f:X\to Y$,
\[
\beta_Y\circ Q_qf=(\mathrm{id}_{C_q}\otimes f)\circ\beta_X.
\]
Writing $\ell_X:S^{0,0}\otimes X\xrightarrow{\sim}X$ for the left
unit isomorphism, the two inclusion compatibilities are
\[
\beta_X\circ\iota_X
  =(\iota_{S^{0,0}}\otimes\mathrm{id}_X)\circ\ell_X^{-1},
\qquad
\varepsilon_{n,X}\circ\iota_{\Sigma^nX}=\Sigma^n\iota_X.
\]
Exactness uses the same suspension comparison $\varepsilon$.
Transport through $\beta$ gives the corresponding shift and exactness
structures on $C_q\otimes-$.
For each $q>0$, the assumptions specify its exact quotient functor
and inclusion-compatible suspension comparison.

In particular, every distinguished triangle
$A\xrightarrow fB\xrightarrow gC\xrightarrow h\Sigma A$ gives a
distinguished row
\[
 Q_qA\xrightarrow{Q_qf}Q_qB\xrightarrow{Q_qg}Q_qC
 \xrightarrow{e_A\circ Q_qh}\Sigma Q_qA,
 \qquad e_A=\varepsilon_{1,A}.
\]
Here $e_A:Q_q(\Sigma A)\xrightarrow{\sim}\Sigma Q_qA$ satisfies
$e_A\circ\iota_{\Sigma A}=\Sigma\iota_A$. The vertices range over all
synthetic objects, for every $q>0$.

\lean{KIP126.Literature.Standard.FiniteQuotientComparison,KIP126.Literature.Standard.finiteQuotientComparison,KIP126.Literature.Standard.finite_quotient_horizontal_row}
\end{definition}
\begin{proof}
The row is the image of the given triangle under the exact functor
$Q_q$, with suspension comparison $e_A=\varepsilon_{1,A}$.
Its inclusion equation follows from the assumed compatibility of
$\varepsilon$ with $\iota$. Thus the same exact functor gives both
the row and its inclusion maps.
\end{proof}

\begin{theorem}[The signed finite quotient row of a normalized triangle]
\label{sem:Literature:normalized-finite-quotient-row}
\level{coarse}
\source{Derived in the fixed synthetic model}
\uses{paper:3.20,sem:AdamsSynthetic:lambda-quotients}
Fix a classical triangle $T$ in the bounded-below finite-type range,
with normalization exponents $e_f+e_g+e_h=1$. Write its normalized
triangle as
\[
 V=(V_1\xrightarrow aV_2\xrightarrow bV_3\xrightarrow c\Sigma V_1),
\]
which is distinguished by Proposition~\ref{paper:3.20}. For $q>0$, write
$Q_qA=\operatorname{cofib}(\lambda_A^q)$ and let $\iota_A:A\to Q_qA$
be the cofiber inclusion. There is an isomorphism
\[
\varepsilon:Q_q(\Sigma V_2)\cong\Sigma Q_qV_2,
\qquad
\varepsilon\circ\iota_{\Sigma V_2}=\Sigma\iota_{V_2},
\]
for which the row
\[
Q_qV_2\xrightarrow{Q_qb}Q_qV_3
\xrightarrow{Q_qc}Q_q(\Sigma V_1)
\xrightarrow{-\varepsilon\circ Q_q(\Sigma a)}
\Sigma Q_qV_2
\]
is distinguished. The construction supplies such an inclusion-compatible $\varepsilon$.
For a given $\varepsilon$, distinguishedness of this row is
equivalent to constructing an isomorphism
\[
e:\operatorname{cofib}(Q_qb)\cong Q_q(\Sigma V_1)
\]
that preserves both structure arrows:
\[
e\circ\iota_{Q_qb}=Q_qc,\qquad
\bigl[-\varepsilon\circ Q_q(\Sigma a)\bigr]\circ e=\delta_{Q_qb}.
\]
Both structure-arrow equations are needed.

\lean{KIP126.Literature.Standard.normalized_finite_quotient_horizontal_row}
\end{theorem}
\begin{proof}
\uses{paper:3.20,sem:Literature:finite-quotient-horizontal-row}
Proposition~\ref{paper:3.20} makes $V$ distinguished. Rotate it to
\[
 V_2\xrightarrow bV_3\xrightarrow c\Sigma V_1
 \xrightarrow{-\Sigma a}\Sigma V_2.
\]
Apply the exact-row conclusion of
Definition~\ref{sem:Literature:finite-quotient-horizontal-row}
to this rotated triangle. It supplies the same inclusion-compatible
suspension isomorphism $\varepsilon$. The exact quotient functor is
additive, so $Q_q(-\Sigma a)=-Q_q(\Sigma a)$; its last arrow is
therefore precisely the displayed signed boundary. This proves the row for all the stated
vertices and every $q>0$.
\end{proof}

\begin{theorem}[The normalized $E_2$ long exact sequence]
\label{sem:Literature:normalized-e2-long-exact}
\level{coarse}
A consequence of the signed quotient row in
\ref{sem:Literature:normalized-finite-quotient-row} at quotient length one.\par
Let the following triangle be distinguished in the bounded-below
finite-type spectrum range:
\[
X\xrightarrow fY\xrightarrow gZ\xrightarrow h\Sigma X,
\qquad e_f=1,\quad e_g=e_h=0.
\]
Its normalized synthetic triangle is distinguished by
Proposition~\ref{paper:3.20}.
Write $Q_1V=V/\lambda$. The normalized map
$\widehat f:\Sigma^{0,1}\nu X\to\nu Y$ induces an additive operation
\[
\partial_f^Q:E_2^{s,t}(X)\longrightarrow E_2^{s+1,t+1}(Y)
\]
through $Q_1(\widehat f)$ and the first-quotient comparisons.
Let $g_2$ be the $E_2$ map of $g$, and let $h_2$ be the
$E_2$ map of $h$ followed by Adams desuspension.
For all $s,t\in\mathbb Z$, the sequence
\[
\cdots\longrightarrow E_2^{s,t+1}(Y)\xrightarrow{g_2}
E_2^{s,t+1}(Z)\xrightarrow{h_2}E_2^{s,t}(X)
\xrightarrow{\partial_f^Q}E_2^{s+1,t+1}(Y)
\xrightarrow{g_2}E_2^{s+1,t+1}(Z)\longrightarrow\cdots
\]
is exact at every displayed position. In particular,
$\ker h_2=\operatorname{im}g_2$,
$\ker\partial_f^Q=\operatorname{im}h_2$, and
$\ker g_2=\operatorname{im}\partial_f^Q$ in the indicated degrees.
\uses{sem:Literature:normalized-finite-quotient-row,sem:AdamsSynthetic:normalized-map,sem:AdamsSynthetic:first-quotient-shifts,sem:AdamsSynthetic:selected-convergence}

\lean{KIP126.Literature.Standard.E2Exactness.e2LongExactSequence}
\end{theorem}
\begin{proof}
\uses{sem:Literature:normalized-finite-quotient-row,sem:AdamsSynthetic:first-quotient-shifts,sem:AdamsSynthetic:selected-convergence}
Write $N=(N_1\to N_2\to N_3\to\Sigma N_1)$ for the normalized
triangle. Apply the signed quotient-row theorem with $q=1$. The
resulting distinguished row is
\[
Q_1N_2\longrightarrow Q_1N_3\longrightarrow Q_1(\Sigma N_1)
\longrightarrow\Sigma Q_1N_2.
\]
The covariant Hom exact sequence gives the kernel--image identities at
all three successive vertices, using rotation for the last one.
First-quotient naturality identifies the first two maps with
$g_2$ and $h_2$. The positive-weight and suspension comparisons
identify the third map with $\partial_f^Q$ after desuspension.
The sign in the quotient row changes neither its kernel nor its image.
Transport along the first-quotient and suspension isomorphisms
preserves these kernel--image identities.
\end{proof}

\section{Products, May push-pull and Toda operations}
The extension and obstruction arguments combine products with cofiber
relations. We first describe multiplication on the positive lambda quotients
and its relation to Ext multiplication. May's tensor construction supplies
the common push-pull witness; Toda witnesses then keep track of the
secondary values and their full indeterminacy. The low-dimensional
identities at the end provide the particular Hopf and theta-five inputs.

\subsection{Products and their detecting labels}

\begin{proposition}[Algebra structures on the positive lambda quotients]
\label{sem:Literature:quotient-algebras}
\level{coarse}
\uses{sem:AdamsSynthetic:lambda-quotients,sem:AdamsSynthetic:graded-hom}
The positive lambda quotients are $E_\infty$ objects by \cite[Remark 3.4]{LWX126}.\par

For each $q>0$, there is a commutative algebra structure on
$S^{0,0}/\lambda^q$ in the symmetric synthetic category whose unit
is the quotient inclusion. Its multiplication and the sphere tensor
comparison induce a bigraded product. Multiplication by a projected
sphere class equals the sphere action. At $q=1$, the next proposition
identifies this product with Ext multiplication.


\lean{KIP126.Literature.Standard.quotientSyntheticSymmetric,KIP126.Literature.Standard.quotientAlgebra,KIP126.Literature.Standard.quotientAlgebra_commutative,KIP126.Literature.Standard.quotientAlgebra_unit,KIP126.Literature.Standard.quotientAlgebraProduct,KIP126.Literature.Standard.quotientAlgebra_sphere_action,KIP126.Literature.Standard.quotientAlgebraConstruction}
\end{proposition}


\begin{proposition}[The first quotient identifies products with Ext products]
\label{sem:Literature:first-quotient-products}
\level{coarse}
\uses{sem:Literature:quotient-algebras,sem:AdamsSynthetic:e2-comparison,sem:AdamsSynthetic:first-quotient-shifts,sem:AdamsSynthetic:sphere-second-page-algebra}
The monoidal comparisons and Ext identification are \cite[Theorems 4.46 and 4.54, Lemma 4.56]{Pst}.\par

For $s,t,s',t'\in\mathbb N$, under the first-quotient equivalence, multiplication of classes in bidegrees $(t-s,t)$ and $(t'-s',t')$ is the classical sphere $E_2$ product in bidegree $(s+s',t+t')$. The action of an untruncated sphere class on a first-quotient class agrees with the same product.


\lean{KIP126.Literature.Standard.sphereFirstQuotient,KIP126.Literature.Standard.quotientAlgebra_first_quotient,KIP126.Literature.Standard.first_quotient_action}
\end{proposition}

\begin{proposition}[Filtration and detection of products]
\label{sem:Literature:filtered-product-detection}
\level{coarse}
\uses{sem:AdamsSynthetic:synthetic-detection-products-exactness,sem:AdamsSynthetic:selected-convergence,sem:Literature:finite-quotient-labels}
For the classical composition pairing, see \cite[Theorem 2.1(iv)]{MossComposition1968}.\par

A sphere class of filtration at least $s$ acting on a class of
filtration at least $s'$ gives a class of filtration at least $s+s'$.
Fix $q>0$ and $s,t,s',t',k,l\in\mathbb N$. Let
$a\in\pi_{t-s,t-k}S^{0,0}$ and
$b\in\pi_{t'-s',t'-l}(S^{0,0}/\lambda^q)$ be detected respectively
by $x\in E_2^{s,t}(S^0)$ and the quotient image of
$y\in E_2^{s',t'}(S^0)$. Then the sphere action $ab$ is detected by
the quotient image of $xy$, in Adams bidegree $(s+s',t+t')$ and
weight $t+t'-k-l$.
The classical product of detected sphere classes is likewise detected
by the product of their classical $E_2$ labels.


\lean{KIP126.Literature.Standard.action_filtration,KIP126.Literature.Standard.finite_action,KIP126.Literature.Standard.classical_product_detection}
\end{proposition}
\begin{proof}
\uses{sem:Literature:first-quotient-products,sem:Literature:infinite-nu-window,sem:Literature:infinite-finite-window,sem:Literature:infinite-shift,sem:Literature:complete-permanent-lifts,sem:AdamsSynthetic:finite-quotient-separation}
For the finite-quotient assertion, first obtain the needed filtration
estimate from the window comparisons. A finite lambda map is surjective
on the relevant associated graded pieces. Lift a filtered class one
piece at a time; the successive corrections terminate at the vanishing
top of the finite filtration. This gives a lift in the same filtration.
Together with the sphere's lambda-divisibility criterion, this writes
the sphere factor as a lambda power and the finite-quotient factor
as the image under a finite lambda map of a diagonal-weight class. The
product of diagonal-weight classes has filtration at least $s+s'$ by
the vanishing for $t<w$. Shifts, lambda multiplication and the
finite lambda maps preserve this estimate.

Write $Q_N=S^{0,0}/\lambda^N$. For $0<r<N$, the actual
$\lambda^r$ action on $Q_N$ factors as the shifted restriction to
$Q_{N-r}$ followed by the finite $\lambda^r$ map. Both maps agree on
the bottom cell; their difference factors through $S^{1,-N-r}$.
The group $\pi_{1,-N-r}Q_N$ vanishes by the finite window and
separation, so the two maps agree.

Suppose $k+l<q$. The windows give
$x\in Z_\infty^{s,t}(S^0)$ and
$y\in Z_{q-l}^{s',t'}(S^0)$. Choose diagonal-weight representatives
$a_0\in\pi_{t-s,t}S^{0,0}$ and
$b_0\in\pi_{t'-s',t'}(S^{0,0}/\lambda^{q-l})$ with these labels.
The first-quotient product comparison identifies the label of
$a_0b_0$ with $xy$: restriction to the first quotient uniquely determines
the second-page label of a diagonal-weight class. Apply the actual $\lambda^k$ action
and the finite $\lambda^l$ map. The window compatibilities carry the
same label to weight $t+t'-k-l$ in the quotient of length $q$.
The resulting factors have the same leading labels as $a,b$; their
differences have filtrations at least $s+1,s'+1$. The preceding
filtration estimate puts the product error in filtration
$s+s'+1$, so the detecting label is unchanged.

If $k+l\ge q$, the product has filtration at least $s+s'$ outside
the finite window, hence is zero by finite-quotient separation; its
projected second-page label is also zero. Thus the argument includes
zero labels and all positive lengths. The classical assertion is the
composition-pairing theorem cited above.
\end{proof}

\subsection{The tensor diagram and its boundary sign}
May's tensor construction gives the following diagram chase, with a
minus sign in the boundary identity.
The unsigned form holds when the homotopy boundary image has exponent two.

\begin{proposition}[Exact tensor functors and May's construction]
\label{sem:Literature:may-tensor}
\level{coarse}
\uses{sem:AdamsSynthetic:categories,sem:AdamsSynthetic:exact-shifts-realization}
Exactness and the suspension comparisons are described in \cite[Proposition 4.2]{Pst} and \cite[Definition 4.1, {\rm(TC1)--(TC2)}]{May01}.\par

Left and right tensoring by a synthetic object are exact: they commute
with suspension and preserve distinguished triangles. We use
suspension isomorphisms compatible with these exact structures and
with the tensor identities in May's construction. The push-pull
assertion below uses May's TC3 compatibility in addition to the
exactness of each tensor functor.


\lean{KIP126.Literature.Standard.tensorLeft_commShift,KIP126.Literature.Standard.tensorRight_commShift,KIP126.Literature.Standard.tensorLeft_isTriangulated,KIP126.Literature.Standard.tensorRight_isTriangulated,KIP126.Literature.Standard.mayData,KIP126.Literature.Standard.mayContextConstruction}
\end{proposition}


% AimPaper: 6.11 lemma 1755 1767
\begin{lemma}[A signed version of Lemma 6.11]
\label{paper:6.11}
\label{sem:Literature:may-push-pull}
\label{lem:452d218c}
\source{AimPaper Lemma 6.11, with the boundary sign corrected}
\level{coarse}
\uses{sem:Literature:may-tensor}
A signed version of \cite{LWX126}, Lemma 6.11, with the tensor and
suspension conventions of \cite{May01}, TC3 and Lemma 4.6. The sign
correction is recorded in the \href{paper/ERRATA.txt}{paper erratum};
compare the \href{paper/main.tex.html#L1755}{original source text}.
\par

Use the exact tensor functors and suspension comparisons of
\ref{sem:Literature:may-tensor}.
Let $T$ and $U$ be distinguished cofiber triangles
\[
 T:X\xrightarrow fY\xrightarrow gZ\longrightarrow\Sus X,
 \qquad
 U:X'\xrightarrow{f'}Y'\xrightarrow{g'}Z'\longrightarrow\Sus X'.
\]
For a synthetic object $A$, write
$\pi_nA=[\Sigma^nS^{0,0},A]=\pi_{n,0}A$ in this lemma.
The corresponding assertions at other weights follow by shifting the
objects and maps.
May's push-pull construction applies to their tensor diagram:
\[
\xymatrix{
    X\sma X' \ar[r] \ar[d] & Y\sma X' \ar[r] \ar[d] & Z\sma X' \ar[d]\\
    X\sma Y' \ar[d] \ar[r] & Y\sma Y' \ar[d] \ar[r] & Z\sma Y' \ar[d]\\
    X\sma Z' \ar[r] & Y\sma Z' \ar[r] & Z\sma Z'\\
}
\]
For every integer $n$, let $a\in\pi_n(X\sma Z')$ and
$b\in\pi_n(Y\sma Y')$ satisfy
\[
 (f\sma1)_*a=(1\sma g')_*b\quad\text{in }\pi_n(Y\sma Z').
\]
There exists a class $c\in\pi_n(Z\sma X')$ satisfying both
\[
 (g\sma1)_*b=(1\sma f')_*c\quad\text{in }\pi_n(Z\sma Y'),
 \qquad
 \partial_{X\sma U}(a)=-\partial_{T\sma X'}(c)
 \quad\text{in }\pi_{n-1}(X\sma X').
\]
Both connecting homomorphisms use these suspension
comparisons. If multiplication by two annihilates the relevant
boundary image, the signed identity also gives the unsigned one.


\lean{KIP126.Literature.Standard.may_pushpull,KIP126.Literature.Standard.may_signed_boundary}
\end{lemma}
\begin{proof}
\uses{sem:Literature:may-tensor}
Use the push-pull object $V$ in Axiom TC3 of \cite{May01}.
Lemma 4.6 supplies a common lift $v$ of $(a,b)$; let $c$ be its
image in $\pi_n(Z\sma X')$. The other tensor square gives the first
image equality.
Desuspending the TC3 boundary relation gives
\[
 \partial_{X\sma U}(a)=-\partial_{T\sma X'}(c).
\]
If the boundary image has exponent two, its negative equals itself,
giving the unsigned form under that hypothesis.
\end{proof}

\subsection{Toda witnesses and the low synthetic values}
A Toda value is represented by cofiber data, and changing the witness
introduces its stated indeterminacy. The synthetic identities below
include the normalization and comparison of the specified Hopf classes.

\begin{definition}[Toda brackets, indeterminacy and naturality]
\label{sem:Literature:toda-relations}
\label{sem:Literature:toda-functoriality}
\level{coarse}
\uses{sem:AdamsSynthetic:categories,sem:AdamsSynthetic:shift-coherence,sem:Literature:may-tensor}
From \cite{ChristensenFrankland2017}, Definition 3.1 and Lemma 3.4.\par

Consider three successive maps
\[
 X\xrightarrow fY\xrightarrow gZ\xrightarrow hW.
\]
We write $\langle f,g,h\rangle$ in the order of these arrows.
A Toda witness for a value $x:\Sigma X\to W$ consists of a
distinguished triangle
\[
 X\xrightarrow fY\xrightarrow iQ\xrightarrow p\Sigma X
\]
and a map $\bar g:Q\to Z$ satisfying
\[
 \bar g\circ i=g,\qquad x\circ p=h\circ\bar g.
\]
Some value admits a witness precisely when the adjacent composites
$g\circ f$ and $h\circ g$ vanish. Fixing one value $x_0$, all values
form the coset
\[
 x_0+\{\,y\circ\Sigma f+h\circ z:
       y:\Sigma Y\to W,\ z:\Sigma X\to Z\,\}.
\]
This is the full Toda indeterminacy. For synthetic sphere classes,
the shifts that make the three maps composable use the bigraded sphere
identifications fixed above.

\emph{Naturality of witnesses.} Pre- and postcomposition carry Toda
witnesses to the corresponding witnesses, and factors can be absorbed
into the first and last entries. Juggling, shuffle and suspension
preserve the cofiber relations. A compatible functor sends a witness to
a witness. Tensoring on either side does the same when equipped with
exactness and shift compatibility. These are the witness laws of
\cite{ChristensenFrankland2017}, Lemma 3.5 and Lemma 5.12.

\lean{KIP126.Literature.Standard.TodaRelation,KIP126.Literature.Standard.TripleToda,KIP126.Literature.Standard.toda_composable,KIP126.Literature.Standard.toda_exists_relation,KIP126.Literature.Standard.toda_relation_iff_indeterminacy,KIP126.Literature.Standard.toda_precompose,KIP126.Literature.Standard.toda_postcompose,KIP126.Literature.Standard.toda_absorb_first,KIP126.Literature.Standard.toda_absorb_last,KIP126.Literature.Standard.toda_juggling,KIP126.Literature.Standard.toda_shuffle_iff,KIP126.Literature.Standard.toda_suspension_iff,KIP126.Literature.Standard.toda_map,KIP126.Literature.Standard.toda_tensor_left,KIP126.Literature.Standard.toda_tensor_right}
\end{definition}
\begin{theorem}[Moss convergence from Massey products to Toda brackets]
\label{sem:Literature:moss-comparison}
\uses{sem:Literature:toda-relations,sem:AdamsSynthetic:mapping-composition}
The convergence theorem is \cite[Theorem 1.1, Definitions 2.3--2.4 and Example 5.1]{BelmontKong2024}.\par

To pass from a Massey product on an Adams page to a Toda value in
homotopy, use the coefficient-unit mapping Adams towers of
Definition~\ref{sem:AdamsSynthetic:mapping-composition}, their composition
pairings and convergence identifications, with each of $W,X,Y,Z$
equal to the completed sphere. The multiplicative tower in this case is
the Adams tower of \cite[Example 5.1]{BelmontKong2024}. Write
$E_r(W,X)$ for its mapping Adams page. For a bidegree $i=(i_1,i_2)$,
put $|i|=i_2-i_1$.
For $r\ge3$, let $a,b,c$ have bidegrees $i,j,k$ on the three
successive mapping pages. Define the page Massey relation using lifts
$a',b',c'$ to $E_{r-1}$ and elements $y,z$ with
$d_{r-1}y=a'b'$ and $d_{r-1}z=b'c'$. The element
$yc'-(-1)^{|i|}a'z$ must represent the Massey member on $E_r$;
use common finite-page representatives for these relations.
Put $d=i+j+k-(r-1,r-2)$.
Suppose $a,b,c$ detect classes $\alpha,\beta,\gamma$,
their two page products vanish, and both stable composites vanish.
Assume weak convergence for the product towers $(W,Y)$ and $(X,Z)$:
writing $T_s$ for the stages of either tower, a class in
$\pi_nT_{s+1}$ which lifts to every $T_t$, $t\ge s+1$,
and maps to zero in $\pi_nT_s$ must be zero.
For each product degree $\ell=i+j,j+k$, exclude every nonzero
$d_m$ on its product tower with source $(q,q+|\ell|+1)$ satisfying
$m>r$, $0\le q<\ell_1-(r-1)$ and $\ell_1<q+m$.
Then some member of the page Massey relation in degree $d$ is permanent
and detects a class $\xi$ whose associated stable map lies in
$\langle\alpha,\beta,\gamma\rangle$.
The detecting page class and the infinite-page class use the same
infinite-cycle representative, allowing zero. The conclusion gives
one permanent Massey member.


\lean{KIP126.Literature.Standard.moss_convergence}
\end{theorem}


\begin{proposition}[The synthetic lift of h-zero]
\label{sem:Literature:synthetic-h-zero}
\level{coarse}
\uses{sem:Literature:hopf-classes,sem:Literature:one-line-differentials,sem:Literature:first-quotient-products}
The normalization is \cite[Example 5.8(4)]{LWX126}; the low-dimensional products are \cite[Proposition A.20 and Remark A.22]{BHS}.\par

There is a class $[h_0]\in\pi_{0,1}S^{0,0}$ with first-quotient label $h_0$ such that $\lambda[h_0]=2$ and $[h_0]\eta=0$. Moreover,
\[[h_0]a=0\qquad\text{for all }a\in\pi_{2,3}S^{0,0}.\]


This lift of $h_0$ will play the role of multiplication by two in the
Toda calculations; the last vanishing controls their indeterminacy.


\lean{KIP126.Literature.Standard.syntheticH0,KIP126.Literature.Standard.syntheticH0_first_quotient_label,KIP126.Literature.Standard.lambdaMultiply_syntheticH0,KIP126.Literature.Standard.sphereProduct_syntheticH0_eta,KIP126.Literature.Standard.sphereProduct_syntheticH0_eq_zero_in_bidegree_two_three,KIP126.Literature.Standard.syntheticH0Construction}
\end{proposition}


\begin{theorem}[The low synthetic Toda value]
\label{sem:Literature:low-toda-comparison}
\level{coarse}
\uses{sem:Literature:toda-relations,sem:Literature:synthetic-h-zero,sem:Literature:hopf-classes}
The Toda value used below appears in \cite[Lemma 7.16, proof]{LWX126}.\par

In the synthetic sphere,
\[\eta^2\in\langle[h_0],\eta,[h_0]\rangle.\]


\lean{KIP126.Literature.Standard.toda_eta_squared}
\end{theorem}

\begin{theorem}[The symmetric Toda value for an order-two class]

\label{sem:Literature:symmetric-toda-comparison}
\level{coarse}
\uses{sem:Literature:toda-relations,sem:Literature:hopf-classes,sem:Literature:synthetic-h-zero}
The classical symmetric Toda law is \cite[Corollary 3.7]{Toda1962};
see also \cite[Theorem 1.3]{BauesMuro2009}. We use the following
synthetic comparison.\par

Assume a stable-enrichment comparison for the synthetic sphere that
sends the classical Hopf class to $\lambda^2\eta$.
For every $\theta\in\pi_{62,64}S^{0,0}$ with $2\theta=0$,
\[\lambda^2\eta\theta\in\langle2,\theta,2\rangle.\]


This assertion applies to every order-two $\theta$, whereas \cite{LWX126}, Lemma 7.16,
uses theta-five. The general classical law is \cite{Toda1962}, Corollary 3.7, and
\cite{BauesMuro2009}, Definition 1.2 (T6) and Theorem 1.3. Its synthetic application
uses this comparison assumption.


\lean{KIP126.Literature.Standard.toda_symmetric}
\end{theorem}


\section{The theta-five survival criterion}
The quadratic construction of Burklund--Xu expresses the obstruction
to the survival of $h_6^2$ as a product involving theta-five. We state
the total-boundary formula first, since it determines which
theta-five class occurs in the criterion.

\subsection{The distinguished theta choice and its total boundary}

\begin{definition}[The distinguished theta-five class and its total boundary]
\label{sem:Literature:bx-choice}
\label{sem:Literature:bx-total-boundary}
\level{coarse}
\uses{sem:Literature:theta-five,sem:Literature:hopf-classes,sem:AdamsSynthetic:lambda-quotients,sem:AdamsSynthetic:standard-classes}
The distinguished class and boundary formula are those of
\cite[Proposition 7.19 and its proof]{BurklundXu}; the present
grading conventions are explained in \cite[Remarks 7.4--7.5]{LWX126}.
Fix the distinguished class
$\theta_5^{\mathrm{BX}}\in\pi_{62,64}S^{0,0}$ with its order-two
relation and first-quotient label $h_5^2$. This synthetic choice is
separate from the classical
choice in Proposition~\ref{sem:Literature:theta-five}. Let $\eta$ be the
normalized synthetic Hopf class with label $h_1$.

Assume that this theta-five class satisfies the following
total-boundary formula.
Let $\delta$ be the connecting homomorphism for
\[
 S^{0,0}\longrightarrow S^{0,0}/\lambda
 \longrightarrow\Sigma^{1,-1}S^{0,0},
\]
with the last suspension desuspended in homotopy. It sends the
first-quotient class $h_6^2\in\pi_{126,128}(S^{0,0}/\lambda)$ to
$\pi_{125,129}S^{0,0}$, and its value is
\[
 \delta(h_6^2)=\lambda\eta(\theta_5^{\mathrm{BX}})^2.
\]
The criterion below uses this $\theta_5^{\mathrm{BX}}$ and this
boundary homomorphism.

\lean{KIP126.Literature.Standard.syntheticTheta5,KIP126.Literature.Standard.syntheticTheta5_detected,KIP126.Literature.Standard.syntheticTheta5_add_self,KIP126.Literature.Standard.eta_first_quotient_label,KIP126.Literature.Standard.etaThetaSquare,KIP126.Literature.Standard.lambdaEtaThetaSquare,KIP126.Literature.Standard.syntheticTheta5Construction,KIP126.Literature.Standard.deltaH6Square,KIP126.Literature.Standard.h6FirstQuotient,KIP126.Literature.Standard.h6FirstQuotient_label,KIP126.Literature.Standard.delta_h6_square_eq_lambda_eta_syntheticTheta5_square}
\end{definition}

\subsection{The reference criterion}

\begin{theorem}[The Burklund--Xu survival criterion]
\label{sem:Literature:bx-survival}
\level{coarse}
\uses{sem:Literature:bx-choice,sem:AdamsSynthetic:lambda-quotients,sem:AdamsSynthetic:standard-classes}
The finite criterion is \cite[Proposition 7.19]{BurklundXu}; its proof
also gives the untruncated criterion.\par

Use the distinguished order-two choice $\theta_5^{\mathrm{BX}}$ above,
as in the quadratic construction of \cite{BurklundXu}. For $r\ge1$, $h_6^2$ survives nontrivially to $E_{r+3}$ if and only if the image of $\eta(\theta_5^{\mathrm{BX}})^2$ in the finite quotient of length $r$ vanishes. In the untruncated sphere, $h_6^2$ is a nonzero permanent cycle if and only if
\[\lambda\eta(\theta_5^{\mathrm{BX}})^2=0.\]


\lean{KIP126.Literature.Standard.h6_square_survives_iff_eta_theta_square_vanishes_mod_lambda,KIP126.Literature.Standard.h6_square_permanent_iff_lambda_eta_theta_square_eq_zero}
\end{theorem}
\begin{proof}
\uses{sem:Literature:bx-total-boundary,sem:Literature:finite-bockstein}
In the quadratic construction of \cite{BurklundXu}, the bottom two
cells identify the total lambda-Bockstein of $h_6^2$ with
$\lambda\eta(\theta_5^{\mathrm{BX}})^2$.
The finite lambda-Bockstein obstruction records the first stage at which
the class fails to lift, giving the displayed criterion in length $r$.
For the untruncated sphere, the same total-boundary formula identifies
permanent lifting with the vanishing of that product. The nonzero
survival assertion also uses that the nonzero second-line class $h_6^2$ has
no incoming Adams boundary.
\end{proof}

Theorem~\ref{paper:7.3} in the Main chapter gives the formulation for
every admissible theta-five choice, including $r=0$, and proves its
equivalence to the existential formulation of \cite{LWX126}, Theorem 7.3.

\section{Framed manifolds and Browder's geometric bridge}
Nonzero survival of a standard square has a geometric meaning through
Browder's criterion: it detects a framed manifold with Kervaire invariant
one. We first describe the manifold and framing, then state the
comparisons needed for this interpretation. The HHR theorem excludes
high dimensions. The Geometry chapter combines these results with the
new permanence result.

\begin{definition}[Framed manifolds and the Kervaire invariant]
\label{sem:Literature:closed-framed-manifolds}
\label{sem:Literature:kervaire-invariant}
\label{sem:Literature:kervaire-existence}
\level{coarse}
Geometric conventions of \cite{LWX126}, Introduction, and
\cite{Browder}, Sections 1--3, especially Theorem 3.2, pp.~166--168.
A closed smooth $n$-manifold is a compact Hausdorff second-countable
smooth manifold without boundary, modelled on $\mathbb R^n$.
A stable tangent framing consists of an integer $r\ge0$ and a smooth
global frame of $TM\oplus\varepsilon^r$: its $n+r$ sections form a
basis at every point. A closed stably framed manifold is the pair
$(M,F)$ of such a manifold and its framing. No connectedness or
unstabilized parallelizability is required.

For a closed stably framed $(4k+2)$-manifold, write
$\operatorname{KI}(M,F)\in\mathbb F_2$ for its Kervaire invariant.
Assume that it agrees with the Arf invariant of the quadratic refinement
of the middle-dimensional mod-two intersection pairing induced by $F$.
Thus the framing is part of its argument. We use this identification
in the geometric statements below; the classical comparison is
\cite[Theorem 3.2]{Browder}.

For $n\in\mathbb N$, write $G(n)$ if there are $k\in\mathbb N$ with
$n=4k+2$ and a closed stably framed $(4k+2)$-manifold $(M,F)$ such that
$\operatorname{KI}(M,F)=1$.

\lean{KIP126.Literature.Geometry.ClosedSmoothManifold,KIP126.Literature.Geometry.StableFraming,KIP126.Literature.Geometry.ClosedStablyFramedManifold,KIP126.Literature.Geometry.kervaireInvariant,KIP126.Literature.Geometry.KervaireInvariantOneExists}
\end{definition}

The invariant in this criterion is the $4k+2$-dimensional invariant of
Definition~\ref{sem:Literature:kervaire-invariant}. Accordingly the square
index is positive, $j\geq1$; the surviving class $h_0^2$ in stem zero is
outside the geometric dimension convention.

% AimPaper: 1.3 theorem 117 123
\begin{definition}[Ordinary classes and their completion]
\label{sem:Literature:pontryagin-thom-completion}
\level{fine}
(For the geometric Pontryagin--Thom interpretation, see \cite{Browder}
and \cite{LWX126}, Theorem 1.3.)\par
\uses{sem:AdamsSynthetic:completion-comparisons}
For every closed stably framed $n$-manifold $M$, $n\in\mathbb N$,
fix a class $\tau(M)\in\pi_n S_{\mathrm{ord}}$, with its framing
understood.
Applying the completion functor and its sphere and suspension
isomorphisms gives $\widehat\tau(M)\in\pi_n S$.
We assume that $\tau(M)$ is identified with the geometric
Pontryagin--Thom class, as in \cite[Section 7]{Browder}. Its completed
class is the image of this ordinary class under completion.

\lean{KIP126.Literature.Geometry.pontryaginThom,KIP126.Literature.Geometry.completedPontryaginThom}
\end{definition}

\begin{theorem}[Browder's criterion through dimension 254]
\label{sem:Literature:finite-browder}
\source{AimPaper Theorem 1.3}
\level{coarse}
This is the finite-range form of \cite[Theorem 7.1, p.~178]{Browder};
see \cite[Theorem 1.3]{LWX126}.\par
Assume the following criterion for $\operatorname{KI}$ and the
standard Adams classes fixed above.
For every integer dimension $0\leq n\leq254$, a closed smooth stably framed manifold of Kervaire
invariant one exists in dimension $n$ if and only if there is an index
$1\leq j\leq7$ such that $n=2^{j+1}-2$ and the square
$h_j^2$ survives with nonzero image at $E_\infty$ in the two-completed sphere
Adams spectral sequence.
\uses{sem:Literature:kervaire-existence,sem:Literature:pontryagin-thom-completion,sem:AdamsSynthetic:standard-classes,sem:SpectralSequence:permanent}

This criterion is assumed under the Arf, Pontryagin--Thom and
classical Adams-label identifications stated above.

\lean{KIP126.Literature.Geometry.browder_iff}
\end{theorem}
\begin{proof}
The classical source is \cite{Browder}, Theorem 7.1, restricted to
$n\leq254$. Since $n=2^{j+1}-2$ and the geometric convention requires
$j\ge1$, this is exactly the range $1\leq j\leq7$.
The theorem relates the geometric Kervaire invariant to nonzero Adams
detection under Pontryagin--Thom. Applying it to the invariant and
Adams classes used here is the comparison assumption in the statement.
\end{proof}

\begin{theorem}[HHR nonexistence in high dimensions]
\label{sem:Literature:hhr-geometric}
\level{coarse}
A geometric consequence of \cite{HHR}, Theorem 1.1, p.~2, and \cite{Browder}, Theorem 7.1.\par
Assume the following nonexistence assertion for the invariant of
Definition~\ref{sem:Literature:kervaire-invariant}.
For every $n\geq254$, no closed smooth stably framed manifold of
Kervaire invariant one exists in dimension $n$.
\uses{sem:Literature:kervaire-existence}

Its geometric interpretation uses the Arf comparison assumed above.

\lean{KIP126.Literature.Geometry.kervaire_invariant_one_not_exists_of_ge_254}
\end{theorem}
\begin{proof}
The high-dimensional nonexistence theorem of \cite{HHR}, combined
with Browder's criterion, has this range under the classical framing
and Kervaire-invariant conventions. Its application to
$\operatorname{KI}$ is the geometric assumption above.
\end{proof}

\section{Classical inputs for the Adams second line}
The second-line classification uses a literature theorem about products
of standard one-line classes, followed by known survival and exclusion
results. Its exhaustiveness input is independent of the bounded source
algebra of Chapter~3. We retain the finite index and degree ranges used
here. The exclusion of $h_7^2$ uses Browder's criterion and HHR
nonexistence, stated above. The corrected differential screening
and the final survivor classification are developed in the later
computation and Main chapters.

\begin{definition}[Products of standard one-line classes]
\label{sem:Literature:two-line-products}
\level{coarse}
From \cite{Adams60}, Theorem 2.5.1(2), p.~45.\par

For $0\leq i,j\leq7$, write
$H_{i,j}=h_i h_j\in E_2^{2,2^i+2^j}(S^0)$ for the product of the
standard one-line classes. Its degree is the sum of their bidegrees.

\lean{KIP126.Literature.Standard.hiProduct}
\end{definition}

\begin{theorem}[Exhaustive standard products on the Adams second line]
\label{sem:E2Computation:two-line-exhaustive}
\level{coarse}
\uses{sem:Literature:two-line-products}
This finite-range statement follows from \cite[Theorem 2.5.1(2), p.~45]{Adams60}.\par

For every integer $t\leq256$ and nonzero $x\in E_2^{2,t}(S^0)$,
there are indices $0\leq i\leq j\leq7$ with
\[
 (1,2^i)+(1,2^j)=(2,t)
\]
such that $x=H_{i,j}$ in $E_2^{2,t}(S^0)$.


\lean{KIP126.Literature.Standard.two_line_exhaustive}
\end{theorem}
\begin{proof}
Apply the cited second-line theorem in the displayed finite internal-degree
range, using the products in \ref{sem:Literature:two-line-products}.
This exhaustiveness is independent of the later differential calculations.
\end{proof}

\begin{theorem}[Classical permanence and exclusions on the second line]

\label{sem:Literature:two-line-classical-inputs}
\level{coarse}
The following assertions concern the $2$-primary sphere Adams spectral
sequence; sources are listed with each assertion.
\begin{enumerate}
\item From \cite{MayThesis1964}, Appendix A; Theorem II.7.9.\par

The three products $H_{0,2},H_{0,3},H_{2,4}$ survive nontrivially to $E_\infty$.
\item From \cite{MayThesis1964}, Appendix A; Theorem II.7.9.\par

For $0\leq j\leq3$, the square $h_j^2$ survives nontrivially.
\item From \cite{MahowaldTangora}, Section 8; \cite{BMMS1986}, VI, Proposition 1.18(ii).\par

The square $h_4^2$ survives nontrivially.
\item From \cite{Mahowaldetaj}, Theorem 1', p.~249.\par

For $3\leq j\leq7$, $H_{1,j}$ survives nontrivially.
\item From \cite{IWX}, Lemma 5.10(1), p.~143.\par

There is a nonzero $d_3$ from $H_{2,5}$ at $(2,36)$ to degree $(5,38)$.
\item From \cite{IWX}, Table 7, p.~225.\par

There is a nonzero $d_4$ from $H_{3,5}$ at $(2,40)$ to degree $(6,43)$.
\item From \cite{IWX}, Table 6, p.~223.\par

There is a nonzero $d_3$ from $H_{3,6}$ at $(2,72)$ to degree $(5,74)$.
\item From \cite{HHR}, Theorem 1.1, p.~2.\par

The square $h_7^2$ does not survive nontrivially to $E_\infty$.
\item From \cite{IWX}, Table 6, p.~222.\par

There is a nonzero $d_3$ from $H_{0,4}$ at $(2,17)$ to degree $(5,19)$.
\end{enumerate}


For clauses 1--2, the $h_0^2$ case follows from Appendix A of \cite{MayThesis1964} and
the grading. For clause 3, a proof is given in \cite{BMMS1986}, VI, Proposition
1.18(ii), p.~183. For clauses 5--7 and 9, \cite[Corollary 1.2]{IWX} transports the
motivic calculations to classical differentials. The nonzero targets
are also displayed in \cite{IWXCharts},
\texttt{Adams-classical-E3.csv}: the source/target row pairs are
$115/111$, $135/130$, $472/457$ and $29/27$, respectively, with the
header counted as the first row.


\lean{KIP126.Literature.Standard.exceptional_hi_products_permanent,KIP126.Literature.Standard.low_hi_squares_permanent,KIP126.Literature.Standard.h4_square_permanent,KIP126.Literature.Standard.h1_hi_product_permanent,KIP126.Literature.Standard.h2_h5_supports_nonzero_d3,KIP126.Literature.Standard.h3_h5_supports_nonzero_d4,KIP126.Literature.Standard.h3_h6_supports_nonzero_d3,KIP126.Literature.Standard.high_hi_squares_not_permanent,KIP126.Literature.Standard.h0_h4_supports_nonzero_d3}
\end{theorem}
\begin{proof}
\uses{sem:Literature:finite-browder,sem:Literature:hhr-geometric}
Use the cited calculations for the low products, squares and
nonzero differentials in the stated finite index ranges.
For $h_7^2$, nonzero survival would imply a framed manifold of Kervaire
invariant one in dimension $254$ by the finite Browder criterion. The
geometric HHR input excludes this dimension.
\end{proof}
